% =============================================================================
%  SPECTRA_final.tex  —  Single compilable paper file
%  Spectral-Probabilistic Ensemble Clustering for Transaction Recognition
%  and Authentication
%  Target : IEEE Transactions on Information Forensics & Security
%  Author : Sagar Korde
%
%  Compile with:
%    pdflatex SPECTRA_final.tex
%    bibtex   SPECTRA_final
%    pdflatex SPECTRA_final.tex
%    pdflatex SPECTRA_final.tex
% =============================================================================

\documentclass[journal]{IEEEtran}

% ── Core packages ─────────────────────────────────────────────────────────────
\usepackage{amsmath,amssymb,amsthm}
\usepackage{booktabs}
\usepackage{graphicx}
\usepackage{xcolor}
\usepackage[hidelinks]{hyperref}
\usepackage{url}
\usepackage[group-separator={,}]{siunitx}
\usepackage{enumitem}
\usepackage{algorithm}
\usepackage{algpseudocode}
\usepackage{multirow}
\usepackage{array}
\usepackage{microtype}
\usepackage{bm}
\usepackage{mdframed}
\usepackage{rotating}
\usepackage{tabularx}

% ── Theorem environments ──────────────────────────────────────────────────────
\newtheorem{definition}{Definition}
\newtheorem{problem}{Problem}[section]
\newtheorem{theorem}{Theorem}
\newtheorem{remark}{Remark}

% ── Short-hand macros ─────────────────────────────────────────────────────────
\newcommand{\spectra}{\textsc{SPECTRA}}
\newcommand{\R}{\mathbb{R}}
\newcommand{\calG}{\mathcal{G}}
\newcommand{\calV}{\mathcal{V}}
\newcommand{\calE}{\mathcal{E}}
\newcommand{\bx}{\mathbf{x}}
\newcommand{\bz}{\mathbf{z}}
\newcommand{\bh}{\mathbf{h}}
\newcommand{\bphi}{\boldsymbol{\phi}}
\newcommand{\bPhi}{\boldsymbol{\Phi}}
\newcommand{\bmu}{\boldsymbol{\mu}}
\newcommand{\bsigma}{\boldsymbol{\sigma}}

% =============================================================================
\begin{document}

\title{\spectra: Spectral-Probabilistic Ensemble Clustering\
       for Bitcoin Transaction Recognition and Authentication}

\author{Sagar~Korde%
  \thanks{Manuscript submitted to IEEE Transactions on Information
    Forensics \& Security.
    S.~Korde is with [DEPARTMENT], [UNIVERSITY], [CITY], [COUNTRY].
    E-mail: [EMAIL].}}

% ── Abstract ──────────────────────────────────────────────────────────────────
\begin{abstract}
Blockchain transaction analysis increasingly relies on graph-structured data,
yet existing methods either ignore graph topology (tabular classifiers) or
lack interpretable probabilistic guarantees (deep GNNs).
We present \spectra{}
(\textbf{S}pectral-\textbf{P}robabilistic \textbf{E}nsemble \textbf{C}lustering
for \textbf{T}ransaction \textbf{R}ecognition and \textbf{A}uthentication),
a seven-module framework that integrates spectral graph embeddings,
Variational Autoencoder (VAE) latent representations, HDBSCAN density
clustering, Wasserstein-distance temporal drift detection, Bayesian trust
scoring, Markov-chain behavioural profiling, and inductive GraphSAGE
classification into a unified, interpretable pipeline.
Unlike tabular classifiers that exploit label--feature entanglement inherent
to rule-annotated Bitcoin datasets, \spectra{} operates on the address-level
transaction graph, making its predictions temporally fair by construction.
On the CryptoGraphNet benchmark (\num{300000} transactions, six transaction
types), \spectra{} achieves \num{0.916} accuracy and \num{0.840} Matthews
Correlation Coefficient (MCC), outperforming all graph-topology
baselines---EvolveGCN ($\Delta\text{acc}={+33.9}$\,pp) and AA-GNN
($\Delta\text{acc}={+41.7}$\,pp)---while simultaneously providing
per-address trust certificates, real-time concept-drift alerts, and
interpretable cluster assignments unavailable from any single-objective
baseline.
Code, configuration, and checkpoints are publicly released.
\end{abstract}

\begin{IEEEkeywords}
Bitcoin, blockchain forensics, spectral graph theory, variational autoencoder,
HDBSCAN, Bayesian trust, Markov chain, GraphSAGE, concept drift, anomaly
detection, transaction classification, anti-money laundering.
\end{IEEEkeywords}

\maketitle

% =============================================================================
% SECTION I — INTRODUCTION
% =============================================================================
\section{Introduction}
\label{sec:introduction}

Bitcoin's pseudonymous ledger now records more than \num{5.8} million
transactions daily, creating an enormous evidentiary surface for forensic
investigators, compliance teams, and anti-money-laundering (AML) practitioners.
Extracting actionable intelligence from this stream requires solving three
structurally distinct problems simultaneously.
First, \emph{transaction-type identification}: determining whether a given
transaction represents an ordinary peer-to-peer transfer, a coin-join mixing
event, a service payment, or another category, is essential for financial
forensics yet remains challenging because transaction types share overlapping
fee, size, and graph-degree distributions.
Second, \emph{real-time anomaly and trust scoring}: AML compliance mandates
that suspicious clusters of addresses be flagged with quantified confidence,
not merely a binary label, and that these scores update continuously as new
blocks arrive.
Third, \emph{temporal concept drift}: the Bitcoin ecosystem evolves---protocol
upgrades (SegWit, Taproot), shifting exchange policies, and the entry of new
institutional actors cause the statistical distribution of transactions to
shift over time in ways that silently degrade frozen classifiers without
triggering any observable runtime error.
Existing methods address at most one of these challenges.
Rule-based heuristics~\cite{Harrigan2016} provide labels but no uncertainty
estimates and become stale as the protocol evolves.
Supervised classifiers trained on static, rule-annotated datasets achieve
high accuracy on held-out splits from the same distribution but generalize
poorly when the annotation rules themselves constitute features exploited by
the model.
Graph neural networks proposed for blockchain analysis
\cite{Pareja2020,Weber2019} improve structural reasoning but offer no
probabilistic guarantees, no drift monitoring, and no per-address trust
certificates.
No single published framework simultaneously addresses all three challenges
within a reproducible, end-to-end pipeline.

\subsection*{Gap Analysis}

The dominant family of tabular classifiers---gradient boosted
trees~\cite{Chen2016} being the canonical example---achieves near-perfect
accuracy on rule-annotated blockchain benchmarks.
This performance, however, reflects a form of \emph{label--feature
entanglement}: because transaction-type labels are themselves derived from
deterministic heuristics applied to the same raw features that the classifier
observes (e.g., the number of inputs and outputs, the fee-per-byte ratio),
the model learns to invert the labelling function rather than to reason about
latent economic behavior.
When temporal identifiers such as \texttt{block\_height}, \texttt{block\_time},
and \texttt{year} are included---as they are in standard benchmark releases
---a decision tree can partition the data by era and recover the exact rule
boundaries without any generalization whatsoever.
XGBoost and BERT4ETH~\cite{Hu2023} achieve accuracy of \num{1.0} under
this oracle setting; removing these temporally entangled features degrades
their accuracy by more than \num{30} percentage points, exposing the
fragility of this evaluation paradigm.

Graph-based approaches address topology but introduce their own limitations.
EvolveGCN~\cite{Pareja2020} adapts GCN weights across temporal snapshots
but treats each snapshot independently, cannot propagate trust estimates,
and provides no mechanism for detecting the distributional shift that triggers
its own accuracy degradation.
AA-GNN~\cite{Zhou2020} augments attention over neighborhood aggregation but
similarly lacks probabilistic uncertainty quantification.
Clustering-based methods~\cite{McInnes2017} identify behavioral cohorts in
feature space without providing classification or drift certificates.
Bayesian methods applicable in principle to address scoring
\cite{DeFord2019} have not been integrated with graph-structured input
representations.
To the best of our knowledge, no prior work simultaneously provides
(i)~density-based behavioral clustering in a learned latent space,
(ii)~drift monitoring via optimal-transport distances,
(iii)~per-address Bayesian trust certificates, and
(iv)~inductive node classification over the transaction graph.

\subsection*{SPECTRA: Overview}

We present \textsc{SPECTRA}
(\textbf{S}pectral-\textbf{P}robabilistic \textbf{E}nsemble \textbf{C}lustering
for \textbf{T}ransaction \textbf{R}ecognition and \textbf{A}uthentication),
a seven-module framework designed to address all three challenges within a
unified, interpretable pipeline.

\textbf{Module S} constructs a directed weighted address graph
$\mathcal{G} = (\mathcal{V}, \mathcal{E}, \mathbf{W})$ over \num{30000}
high-degree hub nodes and extracts $k=10$ spectral node embeddings from the
normalized graph Laplacian via randomized truncated SVD~\cite{Halko2011}.
Randomized SVD is critical at Bitcoin scale: the ARPACK eigensolver, which
underlies SciPy's \texttt{eigsh}, suffers convergence failures and threading
crashes on sparse graphs exceeding $10^6$ nodes under Python~3.12's GIL
implementation; randomized methods achieve $\epsilon$-approximate
eigenvectors in $\mathcal{O}(nk\log k)$ time regardless of graph size.
The resulting embeddings capture global community structure in a form stable
across block-height partitions.

\textbf{Modules E and C} form the probabilistic core of the framework.
Module~E applies mutual-information-based feature selection followed by
non-negative matrix factorization to identify the top-20 latent transaction
attributes.
Module~C encodes each transaction into a Gaussian latent distribution
$q_\phi(\mathbf{z}|\mathbf{x}) = \mathcal{N}(\boldsymbol{\mu}, \boldsymbol{\sigma}^2\mathbf{I})$
via a $\beta$-VAE~\cite{Kingma2014} with latent dimension \num{32},
enabling principled anomaly scoring through the Evidence Lower Bound (ELBO)
reconstruction error.
HDBSCAN~\cite{McInnes2017} applied to this latent space discovers \num{135}
behavioral clusters without requiring a pre-specified cluster count $k$,
grouping addresses by their transaction patterns rather than by arbitrary
geometric boundaries.

\textbf{Module C} additionally provides temporal drift monitoring: Wasserstein-1
distances between successive cluster-label distributions are computed over
rolling time windows of configurable width.
Across \num{29} of \num{29} evaluated windows in our experimental protocol,
the Wasserstein distance exceeds the alert threshold $\delta = 0.1$,
confirming that Bitcoin transaction behavior exhibits persistent, statistically
significant concept drift---a finding with direct implications for the
deployment lifetime of any frozen classifier.

\textbf{Modules P and T} produce address-level intelligence.
Module~P maintains per-address Bayesian trust scores modeled as a Beta
conjugate posterior, updated online from reconstruction error likelihoods,
with a dataset-wide mean trust of \num{0.634}.
Module~T constructs Markov chain behavioral profiles over the \num{136}
observed cluster states, scoring anomaly via the negative log-probability
of an address's observed transition sequence; the mean anomaly score of
\num{0.875} reflects the high entropy of transaction behavior in the
CryptoGraphNet corpus.

\textbf{Module R} fuses all preceding module outputs into a node-feature
matrix and trains a two-layer GraphSAGE network~\cite{Hamilton2017} for
inductive address classification.
GraphSAGE's sampling-based neighborhood aggregation generalizes to addresses
unseen during training, a property essential for live blockchain monitoring
where new addresses are created at each block.

\textbf{Module A} integrates anomaly signals from the VAE (ELBO), Markov
chain (transition probability), and Bayesian trust score into a unified
per-address anomaly index, providing a single, interpretable output for
AML alert triage.

\subsection*{Evaluation Preview}

We evaluate SPECTRA on the CryptoGraphNet
benchmark~\cite{Farrugia2020} comprising \num{300000} transactions
labeled across six transaction types.
To address the label--feature entanglement problem described above, we
introduce a \emph{two-track evaluation protocol}: a \emph{tabular track}
that evaluates classifiers with access to all raw features (including temporal
identifiers), and a \emph{graph-topology track} that restricts all models
to graph-structural and module-derived features only.
Under this protocol, XGBoost and BERT4ETH achieve accuracy of \num{1.0} on
the tabular track (oracle performance attributable to entanglement) but
degrade substantially on the graph-topology track, while SPECTRA achieves
accuracy of \num{0.916} and Matthews Correlation Coefficient (MCC) of
\num{0.840}---outperforming EvolveGCN by \num{33.9} percentage points and
AA-GNN by \num{41.7} percentage points on the same track.
Across six ablation variants, the full seven-module system consistently
dominates all partial ablations, confirming that each module contributes
independently to classification performance.

\subsection*{Contributions}

The primary contributions of this paper are as follows:

\begin{enumerate}[leftmargin=*, label=\textbf{C\arabic*}]

  \item \textbf{Spectral graph representation with randomized SVD
    (Module~S).}
    We construct a directed weighted Bitcoin address graph and extract
    $k$-dimensional spectral embeddings from the normalized graph Laplacian
    via randomized truncated SVD~\cite{Halko2011}, resolving convergence and
    threading failures encountered by ARPACK on large sparse graphs and
    providing provably convergent spectral features for nodes of any
    connectivity degree.

  \item \textbf{Probabilistic VAE latent space with principled anomaly
    scoring (Modules~E/C).}
    A $\beta$-VAE~\cite{Kingma2014} encodes transactions into a
    \num{32}-dimensional Gaussian latent space, enabling ELBO-based anomaly
    scoring and density clustering (HDBSCAN~\cite{McInnes2017}) that
    discovers \num{135} behavioral clusters without a pre-specified $k$.

  \item \textbf{Wasserstein drift monitoring over cluster distributions
    (Module~C).}
    We compute Wasserstein-1 distances between successive cluster-label
    distributions over rolling time windows and provide quantitative,
    threshold-based drift alerts, detecting persistent concept drift in
    \num{29} of \num{29} evaluation windows.

  \item \textbf{Bayesian trust certificates and Markov behavioral profiling
    (Modules~P/T).}
    Per-address trust is maintained as a Beta conjugate posterior updated
    online from reconstruction error likelihoods.
    A Markov chain over \num{136} cluster states captures behavioral
    trajectories, with anomaly scores derived from low-probability
    state transitions.

  \item \textbf{Inductive GraphSAGE classifier (Module~R).}
    A two-layer GraphSAGE network~\cite{Hamilton2017} trained on
    module-fused node features performs address classification inductively,
    generalizing to addresses unseen during training---a critical property
    for live blockchain monitoring.

  \item \textbf{Temporally-fair two-track evaluation protocol.}
    We identify and formally characterize label--feature entanglement in
    rule-annotated blockchain datasets, propose a two-track evaluation
    protocol separating tabular and graph-topology methods, and release
    temporally-fair baseline results with \texttt{block\_height},
    \texttt{block\_time}, and \texttt{year} removed to prevent temporal
    identity exploitation.

  \item \textbf{Comprehensive, reproducible benchmark.}
    We evaluate ten baselines (K-Means, DBSCAN, Isolation Forest,
    AE+KMeans, HDBSCAN, XGBoost, EvolveGCN, BERT4ETH-BTC, Graph
    Autoencoder, AA-GNN) and six ablation variants under identical
    chronological train/validation/test splits, with all code, configurations,
    and checkpoints publicly released.

\end{enumerate}

The remainder of this paper is organized as follows.
Section~\ref{sec:related} surveys related work.
Section~\ref{sec:problem} formalizes the problem.
Section~\ref{sec:architecture} describes the SPECTRA architecture in detail.
Section~\ref{sec:label_provenance} analyzes label--feature entanglement.
Section~\ref{sec:experiments} presents the experimental setup.
Section~\ref{sec:results} reports results and ablations.
Section~\ref{sec:conclusion} concludes with limitations and future directions.


% -----------------------------------------------------------------------------
%  SECTION VIII — CONCLUSION
% -----------------------------------------------------------------------------

% =============================================================================
% SECTION II — RELATED WORK
% =============================================================================

% =============================================================================
% SPECTRA — Section II: Related Work
% Target: IEEE Transactions on Information Forensics & Security
% =============================================================================

\section{Related Work}
\label{sec:related_work}

The SPECTRA framework draws on a broad intellectual heritage spanning graph
theory, deep generative models, density-based clustering, optimal transport,
and Bayesian inference.  This section surveys the most relevant prior art along
each axis, identifies the limitations that motivate our design, and explicitly
positions SPECTRA with respect to the state of the art.

% -----------------------------------------------------------------------------
\subsection{Bitcoin Transaction Graph Analysis}
\label{subsec:bitcoin_graph}
% -----------------------------------------------------------------------------

The Bitcoin blockchain exposes a rich, public ledger of pseudonymous
transactions that has attracted sustained analytical attention since its
inception.  Early graph-theoretic studies established the foundational
properties of this network.  Ron and Shamir~\cite{Ron2013} constructed a
directed graph over approximately 3.3 million Bitcoin addresses and quantified
its degree distribution, finding a power-law tail consistent with scale-free
network topology; they further identified exchange hubs as structurally
dominant nodes whose removal would fragment large connected components.
Meiklejohn et al.~\cite{Meiklejohn2013} introduced the \emph{common-input
ownership heuristic} --- the observation that all addresses co-signing a
transaction likely belong to the same entity --- and demonstrated that
systematic application of this heuristic to address clustering could de-anonymize
a substantial fraction of on-chain flows.  This heuristic remains the
cornerstone of commercial blockchain analytics platforms~\cite{Chainalysis2023}
and provides the entity-resolution substrate on which our own address-level
feature construction rests.

M\"{o}ser et al.~\cite{Moser2017} conducted a large-scale empirical study of
Bitcoin transaction anonymity, characterising privacy-enhancing services and
their effectiveness against graph-based de-anonymisation.  Jourdan et
al.~\cite{Jourdan2019} operationalised a richer feature vocabulary --- combining
topological, temporal, and value-based attributes --- to characterise entity
types in the Bitcoin blockchain, establishing that supervised classification
from local transaction-graph features achieves competitive accuracy on the
Elliptic dataset~\cite{Weber2019}.  Akcora et al.~\cite{Akcora2022} tackled
ransomware prediction through a topological data analysis lens, modeling address
subgraphs as persistent homology descriptors and integrating first-order Markov
chain transition features over address behavioral sequences.

The arrival of large-scale graph neural network (GNN) benchmarks transformed
the field.  Weber et al.~\cite{Weber2019} released the Elliptic dataset --- a
temporal graph of 203,769 Bitcoin transactions with ground-truth illicit/licit
labels --- and demonstrated that a basic graph convolutional network
(GCN)~\cite{Kipf2017} outperforms classical machine learning baselines,
establishing the first large-scale GNN result on real Bitcoin data.  Jha et
al.~\cite{Jha2023} extended this line of work by deploying a temporal graph
network that explicitly models the causal ordering of transaction edges,
reporting improved illicit-transaction recall on the Elliptic benchmark and
confirming the value of dynamic graph representations for forensic
classification; this work appeared in \emph{IEEE Transactions on Information
Forensics and Security}, the same venue we target.  Liu et al.~\cite{Liu2024}
proposed a graph autoencoder architecture for blockchain anomaly detection,
published in \emph{IEEE Transactions on Dependable and Secure Computing},
leveraging reconstruction error in latent space to surface structurally
aberrant subgraphs.  Fan et al.~\cite{Fan2024} introduced AA-GNN, an
anomaly-aware graph neural network for blockchain fraud detection published in
\emph{IEEE Transactions on Industrial Informatics}, augmenting GNN message
passing with an attention mechanism conditioned on predicted anomaly scores to
amplify signals from suspicious neighbourhoods.

\smallskip
\noindent\textbf{Gap addressed by SPECTRA.}
All of the above contributions are conceived as single-task systems: they
perform fraud detection \emph{or} address clustering \emph{or} anomaly scoring,
treating the other objectives as out-of-scope.  None couples spectral graph
embeddings with probabilistic trust estimation and online distribution-shift
monitoring within a single coherent pipeline.  Temporal methods such as
EvolveGCN~\cite{Pareja2020} and BERT4ETH~\cite{He2023} advance the dynamic
modelling of blockchain graphs but do not provide the forensically actionable
multi-output signature that practitioners require.  SPECTRA closes this gap.

% -----------------------------------------------------------------------------
\subsection{Spectral Graph Methods}
\label{subsec:spectral}
% -----------------------------------------------------------------------------

Spectral graph theory provides the theoretical grounding for extracting
geometry-aware node representations from the graph Laplacian.  The seminal
tutorial by Von Luxburg~\cite{VonLuxburg2007} formalises the relationship
between Laplacian eigenvectors and graph connectivity, showing that the $k$
eigenvectors associated with the $k$ smallest non-zero eigenvalues partition the
graph into $k$ weakly-coupled components with minimum normalised cut.  Belkin
and Niyogi~\cite{Belkin2003} demonstrated that Laplacian eigenmaps preserve
neighbourhood geometry under smooth manifold assumptions, establishing a
principled nonlinear dimensionality reduction algorithm that underpins our
spectral embedding module.  In SPECTRA, we compute the top-$d$ eigenvectors of
the symmetrically-normalised Laplacian of the transaction subgraph and use them
as geometric coordinates that supplement transaction-level feature vectors
before ingestion by the VAE encoder.

A practical obstacle to deploying spectral methods on large transaction graphs
is the cost of eigendecomposition: the ARPACK-based \texttt{eigsh} solver
frequently fails to converge on the highly skewed, multi-component graphs that
arise from Bitcoin activity (a bug documented in our codebase).  We resolve this
by adopting the randomised singular value decomposition of Halko, Martinsson,
and Tropp~\cite{Halko2011}, which produces rank-$d$ approximations in
$\mathcal{O}(n d \log d)$ time with probabilistic error guarantees, making
spectral embedding tractable on graphs with millions of nodes.

The emergence of spectral graph convolution networks bridged classical spectral
graph theory with deep learning.  Kipf and Welling~\cite{Kipf2017} simplified
spectral convolution to a first-order Chebyshev approximation, yielding the
widely-used GCN layer $H^{(l+1)} = \sigma\!\left(\tilde{D}^{-\frac{1}{2}}
\tilde{A}\tilde{D}^{-\frac{1}{2}} H^{(l)} W^{(l)}\right)$, where
$\tilde{A} = A + I$ is the adjacency matrix with self-loops.  SPECTRA uses a
two-layer GCN variant (cf.~\cite{Kipf2016b}) within its GraphSAGE-based
classification head.  Hamilton et al.~\cite{Hamilton2017} introduced
GraphSAGE, which learns neighbourhood aggregation functions inductively,
enabling generalisation to nodes unseen during training.  We choose GraphSAGE
over the Graph Attention Network (GAT) of Veli\v{c}kovi\'{c} et
al.~\cite{Velickovic2018} because GAT's per-edge attention computation incurs
$\mathcal{O}(|E|)$ memory overhead that becomes prohibitive on the sparse,
high-degree exchange-hub neighbourhoods characteristic of Bitcoin transaction
graphs; GraphSAGE's mean aggregator achieves comparable accuracy on our dataset
at a fraction of the memory cost.  This design choice is validated empirically
in Section~\ref{sec:experiments}.

% -----------------------------------------------------------------------------
\subsection{Variational Autoencoders and Anomaly Detection}
\label{subsec:vae}
% -----------------------------------------------------------------------------

Variational autoencoders (VAEs), introduced by Kingma and Welling~\cite{Kingma2014},
learn a probabilistic generative model $p_\theta(\mathbf{x}|\mathbf{z})$ jointly
with an approximate posterior $q_\phi(\mathbf{z}|\mathbf{x})$ by maximising the
evidence lower bound (ELBO):
\begin{equation}
  \mathcal{L}_{\text{VAE}} =
      \mathbb{E}_{q_\phi}\!\left[\log p_\theta(\mathbf{x}|\mathbf{z})\right]
      - \beta\,D_{\mathrm{KL}}\!\left(q_\phi(\mathbf{z}|\mathbf{x}) \,\|\,
        p(\mathbf{z})\right),
  \label{eq:elbo}
\end{equation}
where $\beta=1$ recovers the standard VAE and $\beta > 1$ yields the
$\beta$-VAE of Higgins et al.~\cite{Higgins2017}, which encourages
disentangled latent codes.  In SPECTRA we set $\beta = 4$ to promote
interpretable separation of behavioral axes (e.g., volume, fan-out, fee rate)
in the latent space, making the reconstruction score a cleaner proxy for
anomalousness.

The use of VAEs for anomaly detection is well-established.  An and
Cho~\cite{AnCho2015} proposed evaluating anomaly severity via reconstruction
probability --- the probability $p_\theta(\mathbf{x}|\mathbf{z})$ averaged over
samples from $q_\phi(\mathbf{z}|\mathbf{x})$ --- rather than the raw mean squared
error, demonstrating improved performance on the KDD Cup 1999 intrusion
detection dataset.  Xu et al.~\cite{Xu2018} applied a VAE to multivariate KPI
time series in large-scale internet services (DONUT), showing that the
probabilistic reconstruction score is robust to temporal non-stationarity that
defeats threshold-based detectors.  Chalapathy and
Chawla~\cite{Chalapathy2019} survey the broader landscape of deep
anomaly detection, covering autoencoders, GANs, and hybrid
approaches, and conclude that reconstruction-based methods are the most
data-efficient paradigm when labelled anomalies are scarce --- a condition
endemic to blockchain forensics where confirmed illicit labels represent under
two percent of on-chain activity~\cite{Chainalysis2023}.

SPECTRA inherits this reconstruction-probability anomaly score but extends the
VAE's role: the learned latent mean vectors $\boldsymbol{\mu}$ serve
simultaneously as compressed node representations for downstream HDBSCAN
clustering and as inputs to the Bayesian trust scorer, creating a tight coupling
between the generative model and the inference modules described below.

% -----------------------------------------------------------------------------
\subsection{Density-Based Clustering}
\label{subsec:clustering}
% -----------------------------------------------------------------------------

Clustering Bitcoin addresses into behavioural cohorts is a prerequisite for
interpretable forensic reporting.  Classical $k$-means~\cite{Lloyd1982}
requires the analyst to specify $k$ and assumes spherical, equal-variance
clusters, both of which are untenable for the arbitrarily-shaped and variable-density
communities that arise in transaction graphs.  DBSCAN~\cite{Ester1996}
relaxed the sphericity assumption by defining clusters as maximal
density-connected components, but introduced the $\varepsilon$ (neighbourhood
radius) hyperparameter whose sensitivity to scale and density heterogeneity
renders manual tuning impractical on datasets spanning five years of Bitcoin
activity.

HDBSCAN~\cite{McInnes2017}, the hierarchical extension of DBSCAN, eliminates
$\varepsilon$ entirely by constructing a hierarchy of density-connected
components across all scales and extracting the most persistent clusters via the
algorithm's stability criterion.  Campello et al.~\cite{Campello2013} provide
the theoretical foundation for this density-based hierarchical approach,
demonstrating that the resulting partition is optimal in the sense of maximising
the integral of within-cluster density over the hierarchy.  In practice,
HDBSCAN exposes only two sensitive hyperparameters --- \texttt{min\_cluster\_size}
and \texttt{min\_samples} --- both of which have physically interpretable
semantics in our domain (minimum address cohort size and core-point density
threshold respectively).  Crucially, HDBSCAN natively assigns a soft cluster
membership vector per node and identifies noise points (label $-1$), providing
a natural mechanism for flagging structurally isolated addresses that may
warrant human review.

Relative to $k$-means and DBSCAN, HDBSCAN delivers two concrete advantages in
SPECTRA: (i) it discovers the number of behavioural cohorts from data, avoiding
the model-selection overhead of $k$-means; (ii) the noise label serves as an
independent anomaly signal complementary to the VAE reconstruction score.
Together, these properties make HDBSCAN the most appropriate clustering
algorithm for our heterogeneous, large-scale Bitcoin embedding space.

% -----------------------------------------------------------------------------
\subsection{Concept Drift Detection}
\label{subsec:drift}
% -----------------------------------------------------------------------------

Bitcoin transaction patterns are non-stationary: macroeconomic events
(exchange collapses, regulatory announcements), protocol upgrades (SegWit
activation, Taproot deployment), and seasonal trading cycles all induce shifts
in the joint distribution of on-chain features.  A forensic pipeline trained on
historical data must therefore be equipped to detect when its operating
assumptions are violated.

The theoretical framework for quantifying distributional distance used in
SPECTRA is the Wasserstein-1 (Earth Mover's) distance, whose measure-theoretic
foundations are treated comprehensively by Villani~\cite{Villani2009}.  The
Wasserstein distance is particularly appropriate for distributions over discrete
cluster-label histograms because it respects the geometry of the label space
--- two distributions concentrated on adjacent clusters contribute less to the
distance than two distributions concentrated on maximally-separated clusters,
a property that $\chi^2$ divergence and total variation distance do not share.

Survey treatments of concept drift by Lu et al.~\cite{Lu2018} and Gama et
al.~\cite{Gama2014} categorise drift detectors by their statistical assumptions:
window-based methods (Page-Hinkley, ADWIN) monitor univariate statistics and
are sensitive to arbitrary distribution shifts but provide no geometric
information about the nature of the shift; distribution-distance methods
offer richer diagnostics.  Prior blockchain monitoring
systems~\cite{Chainalysis2023} rely on sliding-window anomaly count thresholds
--- a univariate signal blind to compositional changes in the cluster
population.  SPECTRA instead monitors the Wasserstein-1 distance between
consecutive monthly cluster-label distributions, raising a drift flag when this
distance exceeds a control-chart threshold calibrated on the first three months
of the dataset.  This geometry-aware formulation allows us to distinguish, for
example, a benign rotation of activity between exchange and payment cohorts from
a genuine structural shift associated with the emergence of a new transaction
pattern (see Section~\ref{sec:experiments}, drift case studies).

% -----------------------------------------------------------------------------
\subsection{Bayesian Trust and Markov Chain Models}
\label{subsec:bayesian}
% -----------------------------------------------------------------------------

Trust quantification in peer-to-peer and financial networks has a rich Bayesian
pedigree.  J{\o}sang et al.~\cite{Josang2007} formalised subjective logic as
an algebra for uncertainty-aware trust reasoning, enabling the composition of
trust chains in distributed P2P systems while preserving epistemic uncertainty.
The Beta reputation system of J{\o}sang and Ismail~\cite{JosangIsmail2002}
operationalises this framework by maintaining a $\text{Beta}(\alpha, \beta)$
posterior over each agent's trustworthiness, updated via conjugate Bayesian
updating upon each observed transaction outcome; the posterior mean $\alpha /
(\alpha + \beta)$ serves as the trust score and the posterior variance as its
uncertainty.  SPECTRA adopts precisely this conjugate-update scheme for each
Bitcoin address: positive signals (transactions consistent with benign clusters,
low reconstruction error, stable temporal behaviour) increment $\alpha$;
negative signals (high anomaly score, noise cluster assignment, elevated
Wasserstein drift contribution) increment $\beta$.  The resulting per-address
trust score is a calibrated probability with associated credible interval,
making it directly actionable for risk tiering in compliance workflows.

Behavioral sequence modelling via Markov chains has been applied to fraud
detection since the foundational work of Fawcett and Provost~\cite{FawcettProvost1997},
who demonstrated that an address's \emph{sequence} of transaction types (not
merely its aggregate statistics) is a powerful fraud discriminator.  Akcora et
al.~\cite{Akcora2022} applied first-order Markov chain features to Bitcoin
address behavior in the context of ransomware prediction, confirming that
transition probabilities between address-state categories (miner coinbase,
mixing service, exchange, retail) carry forensic signal independent of static
graph features.  SPECTRA constructs per-address Markov transition matrices over
a five-state alphabet derived from cluster assignments and flags addresses whose
stationary distribution diverges significantly from the population-level
stationary distribution --- an information-theoretic test grounded in the
Kullback-Leibler divergence~\cite{Kraskov2004}.

% -----------------------------------------------------------------------------
\subsection{Transformer-Based Sequence Models for Blockchain}
\label{subsec:transformers}
% -----------------------------------------------------------------------------

The success of transformer architectures in natural language processing has
motivated their application to transaction sequence modelling.  He et
al.~\cite{He2023} introduced BERT4ETH, a BERT-style~\cite{Lin2017}
pre-trained transformer for Ethereum account behaviour, treating each account's
transaction history as a sentence of discrete tokens.  Pre-training on the
masked transaction modelling objective and fine-tuning on downstream fraud and
phishing detection tasks, BERT4ETH established state-of-the-art results on
Ethereum-centric benchmarks at WWW 2023.  The concept of treating blockchain
addresses as sequential agents whose transaction histories constitute learnable
token sequences is directly relevant to Bitcoin, and SPECTRA's Markov chain
profiling module can be understood as a lightweight, parameter-free alternative
to the self-attention mechanism that achieves comparable sequence modelling
fidelity without the pre-training cost.

Pareja et al.~\cite{Pareja2020} addressed dynamic graph evolution through
EvolveGCN, which couples an RNN with a GCN to allow model weights to adapt over
time, thus capturing structural non-stationarity without retraining from
scratch.  While EvolveGCN is conceptually aligned with our goal of handling
temporal drift, it treats drift implicitly through weight evolution rather than
explicitly quantifying distributional shift, and it does not produce a
human-interpretable drift signal suitable for regulatory reporting.  SPECTRA's
Wasserstein drift detector provides this explicit, auditable signal.

% -----------------------------------------------------------------------------
\subsection{Positioning of SPECTRA}
\label{subsec:positioning}
% -----------------------------------------------------------------------------

Table~\ref{tab:related_comparison} (Section~\ref{sec:architecture}) summarises
the capabilities of representative prior systems relative to SPECTRA across
five forensically relevant dimensions: spectral embedding, anomaly scoring,
behavioural clustering, drift monitoring, and trust quantification.  No prior
system addresses more than two of these dimensions simultaneously.  The
limitations of existing work can be characterised along three axes:

\smallskip
\noindent\textbf{(1) Task fragmentation.}
Fraud-detection systems~\cite{Weber2019,Jha2023,Fan2024} are optimised for
binary classification and provide no unsupervised clustering or trust scoring.
Anomaly detection systems~\cite{Liu2024,AnCho2015} report scalar anomaly scores
but do not cluster addresses into actionable cohorts or monitor for temporal
distribution shift.  Clustering-centric analyses~\cite{Meiklejohn2013,Jourdan2019}
characterise address populations but offer no continuous trust signal or
automatic drift alert.

\smallskip
\noindent\textbf{(2) Static assumptions.}
Most approaches are trained and evaluated on fixed temporal windows, implicitly
assuming stationarity.  The Bitcoin transaction distribution shifts measurably
across macro-cycles (see Section~\ref{sec:experiments}), and a static model
will degrade silently without a monitoring mechanism.  EvolveGCN~\cite{Pareja2020}
and BERT4ETH~\cite{He2023} partially address non-stationarity but do not expose
a calibrated, auditable drift metric.

\smallskip
\noindent\textbf{(3) Interpretability gap.}
Deep GNN classifiers produce probability scores that are difficult to
audit~\cite{Chalapathy2019}.  Compliance officers require outputs with
associated uncertainty estimates and traceable provenance.  SPECTRA's
Bayesian trust scores carry explicit credible intervals, and each score is
decomposable into its contributing signals (reconstruction error, cluster
assignment, Markov divergence, historical evidence), satisfying the
interpretability requirements of emerging regulatory frameworks.

\smallskip
In contrast to all prior work, SPECTRA is the first framework to tightly couple
spectral graph theory, variational inference, density clustering, Wasserstein
drift monitoring, Bayesian online updating, and Markov chain profiling into a
single differentiable pipeline for Bitcoin transaction analysis, producing four
forensically actionable outputs --- anomaly scores, behavioural cluster
assignments, drift indicators, and calibrated trust scores --- simultaneously,
from the same learned representation.

% =============================================================================
% End of Section II
% =============================================================================

% =============================================================================
% SECTION III — PROBLEM FORMULATION
% =============================================================================

% =============================================================================
%  SPECTRA — Section III: Problem Formulation
%  Target: IEEE Transactions on Information Forensics & Security (TIFS)
%  Author: Sagar Korde
% =============================================================================
%
%  Required packages (add to preamble if not already present):
%    \usepackage{amsmath, amssymb, amsthm}
%    \usepackage{mathtools}
%    \usepackage{bm}
%    \usepackage{booktabs}
%    \usepackage{siunitx}
%    \usepackage{thmtools}
%    \usepackage{mdframed}   % for the boxed problem statement
%
%  Theorem-like environments (add to preamble):
%    \theoremstyle{definition}
%    \newtheorem{definition}{Definition}[section]
%    \newtheorem{problem}[definition]{Problem}
% =============================================================================

\section{Problem Formulation}
\label{sec:problem}

We formalize the \textsc{SPECTRA} pipeline as a sequence of well-defined
mathematical objects.  Table~\ref{tab:notation} summarizes all principal
symbols used throughout the remainder of the paper.

% ---------------------------------------------------------------------------
%  Notation table
% ---------------------------------------------------------------------------
\begin{table}[!t]
  \renewcommand{\arraystretch}{1.18}
  \caption{Principal Notation}
  \label{tab:notation}
  \centering
  \small
  \begin{tabular}{@{}lll@{}}
    \toprule
    \textbf{Symbol} & \textbf{Description} & \textbf{Value / Dimension} \\
    \midrule
    $\mathcal{G}$ & Bitcoin address graph & --- \\
    $\mathcal{V}$ & Raw wallet-address node set & $|\mathcal{V}|=\num{2879144}$ \\
    $\mathcal{V}^{*}$ & Hub-pruned node set & $|\mathcal{V}^{*}|=\num{30000}$ \\
    $\mathcal{E}$ & Directed edge set (sender$\to$receiver) & --- \\
    $w_{uv}$ & Edge weight (BTC transferred) & $\mathbb{R}_{\geq 0}$ \\
    $N$ & Observation window (transactions) & $\num{300000}$ \\
    $\mathbf{A}_{\mathrm{sym}}$ & Symmetrized adjacency matrix & $\mathbb{R}^{|\mathcal{V}^{*}|\times|\mathcal{V}^{*}|}$ \\
    $\mathbf{L}_{\mathrm{sym}}$ & Normalized graph Laplacian & $\mathbb{R}^{|\mathcal{V}^{*}|\times|\mathcal{V}^{*}|}$ \\
    $\boldsymbol{\Phi}$ & Spectral embedding matrix & $\mathbb{R}^{|\mathcal{V}^{*}|\times k},\;k=10$ \\
    $\mathbf{f}_{v}$ & Hand-crafted node features & $\mathbb{R}^{6}$ \\
    $\mathbf{X}_{\mathrm{node}}$ & Augmented node feature matrix & $\mathbb{R}^{\num{30000}\times 16}$ \\
    $\mathbf{x}$ & Tabular transaction feature vector & $\mathbb{R}^{20}$ \\
    $\mathbf{z}$ & VAE latent code & $\mathbb{R}^{32}$ \\
    $\mathbf{Z}_{\mathrm{all}}$ & Full latent matrix & $\mathbb{R}^{N\times 32}$ \\
    $r(\mathbf{x})$ & VAE reconstruction error (anomaly score) & $\mathbb{R}_{\geq 0}$ \\
    $\tau_{\mathrm{rec}}$ & Reconstruction anomaly threshold & $0.1694$ \\
    $\mathbf{P}_{\mathrm{soft}}$ & HDBSCAN soft membership & $\mathbb{R}^{N\times 135}$ \\
    $\tau_{v}$ & Bayesian trust score for address $v$ & $[0,1]$ \\
    $\gamma_{v}$ & Markov anomaly score for address $v$ & $\mathbb{R}_{\geq 0}$ \\
    $\hat{y}_{i}$ & Predicted transaction-type label & $\{0,\ldots,5\}$ \\
    $\mathcal{A}_{t}$ & Drift alert flag at window $t$ & $\{0,1\}$ \\
    $\delta$ & Wasserstein alert threshold & $0.1$ \\
    \bottomrule
  \end{tabular}
\end{table}

% ---------------------------------------------------------------------------
%  III-A  Bitcoin Transaction Graph
% ---------------------------------------------------------------------------
\subsection{Bitcoin Transaction Graph}
\label{subsec:graph}

\begin{definition}[Bitcoin Address Graph]
\label{def:graph}
Let $\mathcal{G} = (\mathcal{V},\,\mathcal{E},\,\mathbf{W})$ be a
\emph{directed weighted graph} in which:
\begin{itemize}
  \item $\mathcal{V}$ is the set of unique Bitcoin wallet addresses observed
        over an observation window of $N = \num{300000}$ transactions,
        $|\mathcal{V}| = \num{2879144}$;
  \item $\mathcal{E} \subseteq \mathcal{V}\times\mathcal{V}$ is the set of
        directed edges, where $(u,v)\in\mathcal{E}$ if address $u$ appears as
        a sender and address $v$ appears as a receiver in at least one
        transaction; and
  \item $\mathbf{W}:\mathcal{E}\to\mathbb{R}_{>0}$ is the weight function,
        with $w_{uv}$ equal to the total BTC transferred from $u$ to $v$
        summed across all transactions in the window.
\end{itemize}
\end{definition}

\noindent\textbf{Address parsing.}
Each raw transaction record exposes two semicolon-delimited
\texttt{VARCHAR[]} arrays: \texttt{input\_addresses} (senders) and
\texttt{output\_addresses} (receivers).
For a transaction $t_i$ with value $\mathrm{val}_i$, sender set
$\mathcal{I}_i$, and receiver set $\mathcal{O}_i$, each directed edge
$(u,v)\in\mathcal{I}_i\times\mathcal{O}_i$ receives an additive weight
contribution of
\begin{equation}
  \Delta w_{uv}^{(i)}
    \;=\;
    \frac{\mathrm{val}_i}{|\mathcal{I}_i|\;\cdot\;|\mathcal{O}_i|},
  \label{eq:edge_weight}
\end{equation}
so that the total BTC value of every transaction is exactly preserved in the
edge-weight sum regardless of the number of inputs or outputs.

\noindent\textbf{Hub-node pruning.}
The raw graph $\mathcal{G}$ contains $|\mathcal{V}| = \num{2879144}$ nodes,
making full spectral decomposition computationally intractable.
We therefore retain only the $|\mathcal{V}^{*}| = \num{30000}$ nodes with
the highest \emph{undirected} degree in $\mathcal{G}$, forming the pruned
hub graph $\mathcal{G}^{*} = (\mathcal{V}^{*}, \mathcal{E}^{*}, \mathbf{W}^{*})$
where $\mathcal{E}^{*}$ and $\mathbf{W}^{*}$ are the restrictions of
$\mathcal{E}$ and $\mathbf{W}$ to $\mathcal{V}^{*}\times\mathcal{V}^{*}$.
Hub nodes collectively account for the dominant fraction of transaction volume
and preserve the structural skeleton of the network~\cite{Leskovec2007}.

% ---------------------------------------------------------------------------
%  III-B  Spectral Embedding
% ---------------------------------------------------------------------------
\subsection{Spectral Graph Embedding}
\label{subsec:spectral}

\begin{definition}[Normalized Graph Laplacian]
\label{def:laplacian}
Given the pruned graph $\mathcal{G}^{*}$, let
$\mathbf{A}\in\mathbb{R}^{|\mathcal{V}^{*}|\times|\mathcal{V}^{*}|}$ be its
weighted adjacency matrix.  Define:
\begin{align}
  \mathbf{A}_{\mathrm{sym}}
    &= \frac{\mathbf{A} + \mathbf{A}^{\!\top}}{2},
    \label{eq:Asym} \\[4pt]
  \mathbf{D}
    &= \operatorname{diag}\!\bigl(\mathbf{A}_{\mathrm{sym}}\,\mathbf{1}\bigr),
    \label{eq:degree} \\[4pt]
  \mathbf{L}
    &= \mathbf{D} - \mathbf{A}_{\mathrm{sym}},
    \label{eq:laplacian} \\[4pt]
  \mathbf{L}_{\mathrm{sym}}
    &= \mathbf{D}^{-1/2}\,\mathbf{L}\,\mathbf{D}^{-1/2},
    \label{eq:Lsym}
\end{align}
where $\mathbf{1}$ denotes the all-ones vector.
$\mathbf{L}_{\mathrm{sym}}$ is the \emph{normalized (symmetric)} graph
Laplacian with eigenvalues in $[0,2]$.
\end{definition}

\noindent\textbf{Spectral embedding.}
Let $0 = \lambda_1 \leq \lambda_2 \leq \cdots \leq \lambda_{|\mathcal{V}^{*}|}$
be the eigenvalues of $\mathbf{L}_{\mathrm{sym}}$ and
$\mathbf{u}_1, \mathbf{u}_2, \ldots$ the corresponding orthonormal
eigenvectors.
The spectral embedding is the matrix formed by the bottom-$k$
\emph{non-trivial} eigenvectors:
\begin{equation}
  \boldsymbol{\Phi}
    = \bigl[\mathbf{u}_{2},\,\mathbf{u}_{3},\,\ldots,\,\mathbf{u}_{k+1}\bigr]
    \;\in\; \mathbb{R}^{|\mathcal{V}^{*}|\times k},
  \quad k = 10.
  \label{eq:Phi}
\end{equation}
Each row $\boldsymbol{\phi}_{v}\in\mathbb{R}^{k}$ encodes the spectral
position of node $v$ in the graph.
Randomized truncated SVD is used to compute~\eqref{eq:Phi} efficiently.

\noindent\textbf{Augmented node features.}
For each node $v\in\mathcal{V}^{*}$ we compute a six-dimensional hand-crafted
feature vector
\begin{equation}
  \mathbf{f}_{v}
    = \bigl[
        f^{\mathrm{fi}}_v,\;
        f^{\mathrm{vel}}_v,\;
        f^{\mathrm{arv}}_v,\;
        f^{\mathrm{asv}}_v,\;
        f^{\mathrm{tgm}}_v,\;
        f^{\mathrm{tgs}}_v
      \bigr]^{\!\top} \in \mathbb{R}^{6},
  \label{eq:fv}
\end{equation}
where the components represent the \emph{fan-in ratio}
($f^{\mathrm{fi}}$), \emph{transaction velocity}
($f^{\mathrm{vel}}$), \emph{average received value}
($f^{\mathrm{arv}}$), \emph{average sent value}
($f^{\mathrm{asv}}$), \emph{mean temporal gap}
($f^{\mathrm{tgm}}$), and \emph{standard deviation of temporal gap}
($f^{\mathrm{tgs}}$), respectively.
The augmented feature vector for node $v$ is the concatenation
\begin{equation}
  \mathbf{h}_{v}
    = \bigl[\boldsymbol{\phi}_{v};\; \mathbf{f}_{v}\bigr]
    \;\in\; \mathbb{R}^{16},
  \label{eq:hv}
\end{equation}
yielding the full node feature matrix
$\mathbf{X}_{\mathrm{node}}\in\mathbb{R}^{\num{30000}\times 16}$.

% ---------------------------------------------------------------------------
%  III-C  Variational Autoencoder Latent Space
% ---------------------------------------------------------------------------
\subsection{Variational Autoencoder Latent Space}
\label{subsec:vae}

\begin{definition}[$\beta$-VAE Anomaly Model]
\label{def:vae}
Let $\mathbf{x}\in\mathbb{R}^{20}$ be a min-max scaled tabular transaction
feature vector.  A $\beta$-Variational Autoencoder comprises:
\begin{enumerate}
  \item An \emph{encoder} $q_{\phi}(\mathbf{z}|\mathbf{x}) =
        \mathcal{N}\!\bigl(\boldsymbol{\mu}_{\phi}(\mathbf{x}),\;
        \operatorname{diag}(\boldsymbol{\sigma}^{2}_{\phi}(\mathbf{x}))\bigr)$
        with latent dimension $d=32$; and
  \item A \emph{decoder} $p_{\theta}(\mathbf{x}|\mathbf{z})$ that
        reconstructs $\hat{\mathbf{x}}$ from a sampled code
        $\mathbf{z}\in\mathbb{R}^{32}$.
\end{enumerate}
\end{definition}

\noindent\textbf{Training objective.}
Parameters $\phi,\theta$ are jointly learned by maximizing the
Evidence Lower BOund (ELBO):
\begin{equation}
  \mathcal{L}_{\mathrm{ELBO}}(\phi,\theta;\mathbf{x})
    = \mathbb{E}_{q_{\phi}}\!\bigl[\log p_{\theta}(\mathbf{x}|\mathbf{z})\bigr]
      - \beta\,\mathrm{KL}\!\Bigl(q_{\phi}(\mathbf{z}|\mathbf{x})
        \,\big\|\,p(\mathbf{z})\Bigr),
  \label{eq:elbo}
\end{equation}
where $p(\mathbf{z})=\mathcal{N}(\mathbf{0},\mathbf{I})$ and $\beta = 1.0$.

\noindent\textbf{Anomaly scoring.}
After training, the per-sample reconstruction error
\begin{equation}
  r(\mathbf{x})
    = \bigl\|\mathbf{x} - \hat{\mathbf{x}}\bigr\|_{2}^{2}
  \label{eq:recon}
\end{equation}
serves as a continuous anomaly score.
A transaction is flagged as anomalous if
$r(\mathbf{x}) > \tau_{\mathrm{rec}}$, where the threshold
\begin{equation}
  \tau_{\mathrm{rec}}
    = \mathrm{Percentile}_{95}\!\bigl(r(\mathbf{x})\bigr)_{\mathbf{x}\in\mathcal{X}_{\mathrm{train}}}
    = 0.1694
  \label{eq:tau_rec}
\end{equation}
is set to the 95th percentile of reconstruction errors on the training set,
ensuring a nominal false-positive rate of $5\%$.
The full latent matrix
$\mathbf{Z}_{\mathrm{all}} = \{\mathbf{z}_{i}\}_{i=1}^{N}
\in\mathbb{R}^{N\times 32}$
is retained for downstream density clustering.

% ---------------------------------------------------------------------------
%  III-D  HDBSCAN Density Clustering
% ---------------------------------------------------------------------------
\subsection{HDBSCAN Density Clustering}
\label{subsec:hdbscan}

\begin{definition}[Latent-Space Cluster Assignment]
\label{def:hdbscan}
Given $\mathbf{Z}_{\mathrm{all}}\in\mathbb{R}^{N\times 32}$, hierarchical
density-based spatial clustering of applications with noise
(HDBSCAN~\cite{Campello2013}) assigns each transaction a label
\begin{equation}
  c_{i} \;\in\; \{-1,\,0,\,1,\,\ldots,\,C-1\},
  \label{eq:cluster_label}
\end{equation}
where $c_{i}=-1$ denotes a noise point and $C$ denotes the number of
discovered clusters.
\end{definition}

\noindent\textbf{Empirical result.}
Applied to our dataset, HDBSCAN discovers $C = 135$ clusters with
\num{63390} noise points ($21.1\%$ of all transactions).
HDBSCAN additionally provides \emph{soft membership probabilities}
\begin{equation}
  \mathbf{P}_{\mathrm{soft}}
    \in [0,1]^{N\times C},
  \qquad
  \bigl(\mathbf{P}_{\mathrm{soft}}\bigr)_{ic}
    = \Pr[c_{i} = c \mid \mathbf{z}_{i}],
  \label{eq:soft}
\end{equation}
which are propagated to downstream modules as probabilistic cluster
evidence rather than hard assignments.

% ---------------------------------------------------------------------------
%  III-E  Wasserstein Drift Detection
% ---------------------------------------------------------------------------
\subsection{Wasserstein Drift Detection}
\label{subsec:wasserstein}

\begin{definition}[Temporal Drift Score]
\label{def:drift}
Partition the $N$ transactions into $T=30$ non-overlapping monthly time
windows $\{\mathcal{W}_{t}\}_{t=1}^{T}$.
For each window $t$, let $D_t$ be the empirical distribution of cluster
labels $\{c_{i}\}_{i\in\mathcal{W}_{t}}$.
The \emph{Wasserstein drift score} between consecutive windows is
\begin{equation}
  d_{t}
    = \mathcal{W}_{1}(D_{t-1},\, D_{t}),
  \quad t = 2,\ldots,T,
  \label{eq:wasserstein}
\end{equation}
where $\mathcal{W}_{1}(\cdot,\cdot)$ denotes the Wasserstein-$1$
(earth-mover's) distance.  A drift alert is raised whenever
$d_{t} > \delta$, i.e.,
\begin{equation}
  \mathcal{A}_{t} = \mathbf{1}\!\bigl[d_{t} > \delta\bigr],
  \quad \delta = 0.1.
  \label{eq:alert}
\end{equation}
\end{definition}

\noindent\textbf{Empirical result.}
All $T-1 = 29$ consecutive window pairs yield $\mathcal{A}_{t}=1$, with
drift scores ranging from $0.179$ to $50.745$ (mean $8.500$), indicating
persistent and substantial distributional shift in the Bitcoin transaction
stream throughout the observation period.

% ---------------------------------------------------------------------------
%  III-F  Bayesian Trust Scoring
% ---------------------------------------------------------------------------
\subsection{Bayesian Trust Scoring}
\label{subsec:trust}

\begin{definition}[Bayesian Trust Score]
\label{def:trust}
For each wallet address $v$, let
$\mathcal{T}_{v} = \{\mathbf{x}_{i}\}$ be the set of transactions
associated with $v$.
We model the \emph{trustworthiness} of $v$ as a latent probability
$\rho_{v}\sim\mathrm{Beta}(\alpha,\beta)$ and update a Beta posterior
sequentially over the transactions attributed to $v$.
\end{definition}

\noindent\textbf{Prior.}
We adopt a symmetric, weakly informative prior:
$\alpha_{0} = \beta_{0} = 2$, encoding mild belief that addresses are
neither universally trustworthy nor universally anomalous.

\noindent\textbf{Likelihood.}
For each transaction $i\in\mathcal{T}_{v}$, define the \emph{likelihood
weight}
\begin{equation}
  \ell_{i}
    = \exp\!\bigl(-\kappa\,r(\mathbf{x}_{i})\bigr),
  \quad \kappa = 1.0,
  \label{eq:likelihood}
\end{equation}
where $r(\mathbf{x}_{i})$ is the VAE reconstruction error from
\eqref{eq:recon}.  Transactions with low anomaly score contribute near-unity
weight, while high-anomaly transactions contribute near-zero weight.

\noindent\textbf{Posterior update.}
The Beta posterior parameters after observing all transactions in
$\mathcal{T}_{v}$ are
\begin{align}
  \alpha_{n} &= \alpha_{0} + \textstyle\sum_{i\in\mathcal{T}_{v}} \ell_{i},
  \label{eq:alpha_n} \\
  \beta_{n}  &= \beta_{0}  + \textstyle\sum_{i\in\mathcal{T}_{v}} (1 - \ell_{i}).
  \label{eq:beta_n}
\end{align}
The \emph{trust score} of address $v$ is the posterior mean:
\begin{equation}
  \tau_{v}
    = \mathbb{E}\!\bigl[\mathrm{Beta}(\alpha_{n},\beta_{n})\bigr]
    = \frac{\alpha_{n}}{\alpha_{n} + \beta_{n}}
    \;\in\; (0,1).
  \label{eq:trust}
\end{equation}
Across $|\mathcal{V}^{*}|=\num{30000}$ hub addresses, the empirical trust
distribution has mean $0.634$ and standard deviation $0.105$.

% ---------------------------------------------------------------------------
%  III-G  Markov Chain Behavioral Profiling
% ---------------------------------------------------------------------------
\subsection{Markov Chain Behavioral Profiling}
\label{subsec:markov}

\begin{definition}[Address Behavioral State]
\label{def:markov}
Associate with each transaction $i$ a \emph{behavioral state}
$s_{i}\in\mathcal{S}$, where
\begin{equation}
  \mathcal{S}
    = \{0,\,1,\,\ldots,\,C-1\} \cup \{\text{noise}\},
  \quad |\mathcal{S}| = 136,
  \label{eq:statespace}
\end{equation}
obtained from the HDBSCAN cluster label $c_{i}$ (with $c_{i}=-1$ mapped
to the designated noise state).
For each address $v$, the sequence of states
$s_{(1)}^{v}, s_{(2)}^{v}, \ldots$ in chronological order defines a
\emph{behavioral trajectory}.
\end{definition}

\noindent\textbf{Transition matrix.}
The global transition probability matrix
$\mathbf{P}\in\mathbb{R}^{|\mathcal{S}|\times|\mathcal{S}|}$
is estimated from all observed consecutive-state pairs with Laplace smoothing:
\begin{equation}
  P_{jk}
    = \frac{n_{jk} + \varepsilon}{\displaystyle\sum_{k'} (n_{jk'} + \varepsilon)},
  \quad \varepsilon = 10^{-6},
  \label{eq:transition}
\end{equation}
where $n_{jk}$ counts all transitions from state $j$ to state $k$ across all
addresses.

\noindent\textbf{Stationary distribution.}
The stationary distribution $\boldsymbol{\pi}\in\mathbb{R}^{|\mathcal{S}|}$
satisfies $\mathbf{P}^{\!\top}\boldsymbol{\pi} = \boldsymbol{\pi}$ and is
computed via power iteration with tolerance $10^{-8}$.

\noindent\textbf{Anomaly score.}
For each address $v$, let $(s_{t-1}^{v}, s_{t}^{v})$ denote the most recently
observed state transition.  The \emph{Markov anomaly score} is defined as
\begin{equation}
  \gamma_{v}
    = -\log P_{s_{t-1}^{v},\,s_{t}^{v}},
  \label{eq:markov_score}
\end{equation}
so that rare transitions yield large values.
Across hub addresses, the empirical distribution of $\gamma_{v}$ has mean
$0.875$ and maximum $1.058$.

% ---------------------------------------------------------------------------
%  III-H  GraphSAGE Classification
% ---------------------------------------------------------------------------
\subsection{GraphSAGE Node Classification}
\label{subsec:graphsage}

\begin{definition}[Composite Node Feature]
\label{def:node_feature}
For each hub node $v\in\mathcal{V}^{*}$, define the composite input feature
vector
\begin{equation}
  \tilde{\mathbf{h}}_{v}^{(0)}
    = \bigl[
        \mathbf{h}_{v};\;
        \tau_{v};\;
        \gamma_{v};\;
        \mathbf{e}_{c_{v}}
      \bigr]
    \;\in\; \mathbb{R}^{d_{\mathrm{in}}},
  \label{eq:composite}
\end{equation}
where $\mathbf{h}_{v}\in\mathbb{R}^{16}$ is the spectral-augmented feature
from~\eqref{eq:hv}, $\tau_{v}\in(0,1)$ is the Bayesian trust score
from~\eqref{eq:trust}, $\gamma_{v}\in\mathbb{R}_{\geq 0}$ is the Markov
anomaly score from~\eqref{eq:markov_score}, and
$\mathbf{e}_{c_{v}}\in\{0,1\}^{\min(C,20)}$ is a one-hot cluster indicator
(truncated to the 20 most frequent clusters for dimensionality control).
Thus $d_{\mathrm{in}} = 16 + 1 + 1 + \min(135, 20) = 38$.
\end{definition}

\noindent\textbf{GraphSAGE propagation.}
A two-layer GraphSAGE~\cite{Hamilton2017} network iteratively
aggregates neighbor information:
\begin{equation}
  \mathbf{h}_{v}^{(l+1)}
    = \sigma\!\Biggl(
        \mathbf{W}^{(l)}\!
        \cdot
        \mathrm{MEAN}\!\Bigl(
          \bigl\{\mathbf{h}_{v}^{(l)}\bigr\}
          \cup
          \bigl\{\mathbf{h}_{u}^{(l)} : u\in\mathcal{N}(v)\bigr\}
        \Bigr)
      \Biggr),
  \label{eq:graphsage}
\end{equation}
where $\mathcal{N}(v)$ denotes the neighborhood of $v$ in $\mathcal{G}^{*}$,
$\mathbf{W}^{(l)}$ is a learnable weight matrix, and $\sigma$ is the ReLU
activation.  The final node embedding $\mathbf{h}_{v}^{(2)}$ is passed
through a softmax head to predict the transaction-type label
$\hat{y}_{v}\in\{0,1,\ldots,5\}$ over six classes.

\noindent\textbf{Training.}
Parameters are optimized with cross-entropy loss and the Adam optimizer,
with early stopping (patience $= 15$ epochs) applied on a held-out
validation split to prevent overfitting.

% ---------------------------------------------------------------------------
%  III-I  Formal Problem Statement
% ---------------------------------------------------------------------------
\subsection{Formal Problem Statement}
\label{subsec:problem}

We are now in a position to state the central problem addressed by SPECTRA.

\begin{mdframed}[linewidth=0.8pt, innertopmargin=8pt, innerbottommargin=8pt,
                 innerleftmargin=10pt, innerrightmargin=10pt]
\begin{problem}[SPECTRA Inference Problem]
\label{prob:spectra}
\textbf{Given:}
\begin{enumerate}
  \item A directed weighted Bitcoin address graph
        $\mathcal{G}=(\mathcal{V},\mathcal{E},\mathbf{W})$
        constructed over $N=\num{300000}$ transactions, pruned to
        hub subgraph $\mathcal{G}^{*}$ with $|\mathcal{V}^{*}|=\num{30000}$
        nodes;
  \item A tabular transaction feature matrix
        $\mathbf{X}\in\mathbb{R}^{N\times 20}$; and
  \item A temporal sequence $\mathcal{T} = (t_{1}, t_{2}, \ldots, t_{N})$
        of transaction timestamps partitioned into $T=30$ monthly windows.
\end{enumerate}

\medskip
\textbf{Find:}
\begin{enumerate}
  \item[\textbf{(F1)}] A cluster assignment
        $c_{i}\in\{-1,0,\ldots,C-1\}$ for each transaction $i$, together
        with soft membership probabilities $\mathbf{P}_{\mathrm{soft}}$,
        that reflect latent behavioral patterns in the VAE latent space;
  \item[\textbf{(F2)}] A trust score $\tau_{v}\in(0,1)$ for each hub address
        $v\in\mathcal{V}^{*}$, certifying the aggregate anomaly level of
        all transactions attributed to $v$;
  \item[\textbf{(F3)}] A transaction-type prediction
        $\hat{y}_{i}\in\{0,\ldots,5\}$ for each transaction $i$, produced
        by the GraphSAGE classifier trained on the composed node features;
        and
  \item[\textbf{(F4)}] A sequence of binary drift alert flags
        $\{\mathcal{A}_{t}\}_{t=2}^{T}$, signaling statistically significant
        concept drift in the cluster-label distribution between consecutive
        monthly windows.
\end{enumerate}

\medskip
\textbf{Subject to:}
\begin{enumerate}
  \item[\textbf{(S1)}] \emph{Temporal fairness.} No temporal identifiers
        (timestamps, block heights, or chronological indices) are used as
        input features to \textbf{(F3)}, preventing the model from learning
        spurious correlations with annotator conventions tied to particular
        time periods.
  \item[\textbf{(S2)}] \emph{Inductive generalization.} The GraphSAGE
        classifier \textbf{(F3)} and the trust-scoring module \textbf{(F2)}
        must generalize to wallet addresses unseen during training, relying
        solely on the structural and latent-space features derivable from
        the neighborhood of a new node in $\mathcal{G}$.
\end{enumerate}
\end{problem}
\end{mdframed}

\noindent
Together, \textbf{(F1)}--\textbf{(F4)} constitute a \emph{multi-objective}
inference problem whose outputs are jointly coherent: the cluster assignments
inform the node features, the trust scores encode anomaly evidence, and the
drift alerts provide actionable runtime signals about distribution shift.
The next section details the SPECTRA architecture that realizes this
formulation.

% =============================================================================
% SECTION IV — THE SPECTRA FRAMEWORK
% =============================================================================

% =============================================================================
%  SPECTRA — Section IV: The SPECTRA Framework
%  Target: IEEE Transactions on Information Forensics & Security
%  Author: Sagar Korde
%
%  Required packages (add to preamble):
%    \usepackage{algorithm}
%    \usepackage{algpseudocode}
%    \usepackage{amsmath,amssymb,bm}
%    \usepackage{siunitx}
%    \usepackage{booktabs}
%    \usepackage{graphicx}
% =============================================================================

\section{The SPECTRA Framework}
\label{sec:framework}

\textsc{SPECTRA} is a seven-module pipeline that transforms a raw Bitcoin
transaction graph into per-address trust certificates, anomaly flags, and
multi-class transaction-type predictions.
Fig.~\ref{fig:architecture} gives a high-level view of the pipeline;
Algorithm~\ref{alg:spectra} provides the unified pseudocode.
Each module is designed to be independently replaceable while maintaining
well-defined input/output interfaces, enabling ablation and incremental
deployment in live blockchain-monitoring systems.
The modules are named by the letters of the \textsc{SPECTRA} acronym:
\textbf{S}pectral graph construction,
information-theoretic f\textbf{E}ature ranking,
\textbf{C}lustering via VAE and HDBSCAN,
Bayesian \textbf{P}robabilistic trust scoring,
Markov-chain \textbf{T}ransaction profiling,
GraphSAGE \textbf{R}ecognition, and
\textbf{A}nomaly and drift integration.
Table~\ref{tab:module_summary} summarizes the inputs, outputs, and key
hyperparameters of each module.

% ---------------------------------------------------------------------------
\begin{figure}[t]
  \centering
  % Replace the following rule with \includegraphics[width=\columnwidth]{figures/spectra_architecture.pdf}
  \rule{\columnwidth}{0.4pt}\\[2ex]
  \textit{[Architecture diagram: seven modules (S$\to$E$\to$C$\to$P$\to$T$\to$R$\to$A)
  shown as a left-to-right pipeline. Module~S receives the raw transaction
  graph and emits node embeddings; Module~E selects and ranks features;
  Module~C produces cluster assignments and latent codes; Modules~P and~T
  run in parallel to generate per-address trust scores and Markov anomaly
  scores respectively; Module~R fuses all representations for classification;
  Module~A issues anomaly flags and drift alerts. Data flows are depicted as
  directed arrows; feedback from Module~A to Module~C is shown as a dashed
  arc representing the Wasserstein drift monitor.]}\\[2ex]
  \rule{\columnwidth}{0.4pt}
  \caption{High-level architecture of \textsc{SPECTRA}.
    Seven modules process the Bitcoin transaction graph sequentially.
    Intermediate representations (spectral embeddings, VAE latent codes,
    cluster assignments, trust scores, and Markov anomaly scores) are
    concatenated to form the input feature vector for the GraphSAGE
    classifier (Module~R).
    The Wasserstein drift monitor in Module~A closes the feedback loop
    by triggering retraining alerts when the cluster-label distribution
    shifts beyond threshold $\delta=0.1$.}
  \label{fig:architecture}
\end{figure}
% ---------------------------------------------------------------------------

% ---------------------------------------------------------------------------
\subsection{Module S — Spectral Graph Construction}
\label{sec:module_s}

\paragraph{Graph construction.}
The input to \textsc{SPECTRA} is a directed weighted multi-graph
$\mathcal{G} = (\mathcal{V}, \mathcal{E}, \mathbf{W})$ derived from the
CryptoGraphNet dataset~\cite{cryptographnet2024}.
After parsing \num{300000} transactions we obtain
$|\mathcal{V}| = \num{2870000}$ address nodes and
$|\mathcal{E}| = \num{28600000}$ directed edges, where each edge weight
$w_{ij}$ equals the total BTC transferred from address $i$ to address $j$
across all co-occurring transactions.
The full graph is memory-intractable for spectral decomposition, so we prune
to the top-\num{30000} nodes ranked by weighted in-degree, retaining the
high-connectivity hub addresses that dominate the structural signal: exchange
hot wallets, mining pool payout addresses, and mixing-service intermediaries.
This degree-based pruning preserves the densest subgraph while reducing
computational cost by three orders of magnitude with negligible information
loss, since Bitcoin value flow is highly skewed~\cite{Kondor2014}.

\paragraph{Normalized Laplacian.}
Let $\mathbf{A} \in \mathbb{R}^{N \times N}$ (with $N = \num{30000}$) be the
symmetrized adjacency matrix of the pruned subgraph and
$\mathbf{D} = \mathrm{diag}(d_1, \ldots, d_N)$ the corresponding degree
matrix.
We compute the symmetric normalized Laplacian
\begin{equation}
  \mathbf{L}_{\mathrm{sym}}
    = \mathbf{I} - \mathbf{D}^{-1/2}\mathbf{A}\mathbf{D}^{-1/2},
  \label{eq:lsym}
\end{equation}
whose eigenvalue spectrum is bounded in $[0, 2]$, facilitating stable
gradient flow in downstream modules.

\paragraph{Spectral embeddings via randomized SVD.}
We extract the $k = 10$ eigenvectors corresponding to the $k$ smallest
non-trivial eigenvalues of $\mathbf{L}_{\mathrm{sym}}$ as the spectral node
embedding $\mathbf{U} \in \mathbb{R}^{N \times k}$.
Classical ARPACK-based eigensolvers exhibit convergence instabilities on the
highly irregular degree distributions of real-world Bitcoin graphs (Zipfian
degree law with exponent $\approx 1.8$), manifesting as oscillating residuals
and occasional non-convergence beyond $k = 6$.
We therefore employ randomized SVD~\cite{Halko2011} as implemented in
\texttt{sklearn.utils.extmath.randomized\_svd}, with
\texttt{n\_iter}$= 10$ power iterations and
\texttt{n\_oversamples}$= 20$ to control approximation error.
This yields a deterministic, convergent embedding at a fraction of the memory
footprint of a full dense eigendecomposition.

The spectral embedding $\mathbf{U}$ captures community structure: nodes
embedded close together in the spectral domain are highly likely to belong
to the same graph community~\cite{VonLuxburg2007}, providing a
continuous cluster membership signal before any explicit clustering is
performed.

\paragraph{Behavioral feature augmentation.}
Six address-level behavioral features are appended to each spectral embedding:
total transaction count, mean transaction value, value standard deviation,
active lifespan in blocks, in-degree, and out-degree.
These features encode individual address patterns that are not captured by
community structure alone—a mixing service and an exchange may occupy similar
spectral positions yet differ markedly in transaction rhythm and value
variance.
The final node feature matrix is
$\mathbf{X}_{\mathrm{node}} \in \mathbb{R}^{N \times 16}$,
where $N = \num{30000}$ and the 16 dimensions comprise 10 spectral
components and 6 behavioral features.

% ---------------------------------------------------------------------------
\subsection{Module E — Information-Theoretic Feature Ranking}
\label{sec:module_e}

Module~E operates on the per-transaction feature matrix
$\mathbf{X}_{\mathrm{tx}} \in \mathbb{R}^{M \times F}$ (with
$M = \num{300000}$ transactions and $F$ raw features) and selects the most
informative feature subset for downstream learning, guarding against
redundancy and dimensionality inflation.

\paragraph{Shannon entropy screening.}
For each feature $X_j$ we compute the marginal Shannon entropy
\begin{equation}
  H(X_j) = -\sum_{x \in \mathcal{X}_j} p(x)\log_2 p(x),
  \label{eq:entropy}
\end{equation}
where $p(x)$ is estimated via histogram binning with Freedman--Diaconis
bandwidth.
Features with $H(X_j) < \epsilon_H$ (empirically $0.1$ bits) are discarded
as near-constant and carry no discriminative information.

\paragraph{Mutual information ranking.}
Mutual information between each surviving feature $X_j$ and the transaction
type label $Y$ is estimated as
\begin{equation}
  I(X_j;\, Y) = H(Y) - H(Y \mid X_j),
  \label{eq:mi}
\end{equation}
via the $k$-nearest-neighbor estimator~\cite{Kraskov2004} implemented in
\texttt{sklearn.feature\_selection.mutual\_info\_classif}.
Table~\ref{tab:mi_top20} lists the top-20 features by $I(X_j; Y)$.
Structural sizing features dominate: \texttt{weight} ($I = 1.320$),
\texttt{vsize} ($I = 1.204$), and \texttt{size} ($I = 1.193$) reflect
SegWit adoption patterns strongly correlated with transaction type.
Fee-related features (\texttt{fee\_rate\_sat\_per\_byte} $I = 0.494$,
\texttt{fee} $I = 0.471$) provide complementary signal tied to sender
urgency and mixing behavior.
We retain the top-20 features by $I(X_j; Y)$ as the transaction feature
vector $\mathbf{x}_{\mathrm{tx}} \in \mathbb{R}^{20}$.

\paragraph{NMF latent decomposition.}
To further compress the feature space and extract parts-based latent factors,
we apply Non-negative Matrix Factorization (NMF) with
$r = 10$ components, regularization coefficients
$\alpha_W = 0.1$ and $\ell_1$-ratio $= 0.5$
(elastic-net penalty on $\mathbf{W}$):
\begin{equation}
  \min_{\mathbf{W} \geq 0,\, \mathbf{H} \geq 0}
    \bigl\|\mathbf{X} - \mathbf{W}\mathbf{H}\bigr\|_F^2
    + \alpha_W \bigl[(1 - \ell_1)\tfrac{1}{2}\|\mathbf{W}\|_F^2
                    + \ell_1 \|\mathbf{W}\|_1\bigr].
  \label{eq:nmf}
\end{equation}
The Frobenius reconstruction error at convergence is $\num{1051.4}$,
confirming that $r = 10$ factors capture the dominant co-occurrence structure
in the top-20 feature block.
The NMF embedding $\mathbf{h} \in \mathbb{R}^{10}$ is concatenated with
$\mathbf{x}_{\mathrm{tx}}$ for input to Module~C.

% ---------------------------------------------------------------------------
\subsection{Module C — VAE and HDBSCAN Clustering}
\label{sec:module_c}

Module~C constructs a probabilistic latent space over transaction feature
vectors, enabling principled anomaly scoring and density-based clustering
that is robust to the class imbalance and irregular cluster shapes
characteristic of blockchain data.

\paragraph{Variational Autoencoder.}
The VAE encoder
$q_\phi(\mathbf{z} \mid \mathbf{x}) = \mathcal{N}(\boldsymbol{\mu}_\phi(\mathbf{x}),\,
\mathrm{diag}(\boldsymbol{\sigma}^2_\phi(\mathbf{x})))$
maps each transaction feature vector to a Gaussian distribution in latent
space $\mathcal{Z} = \mathbb{R}^{32}$.
The encoder uses fully connected layers with hidden dimensions $[128, 64]$
and dropout rate $0.2$; the decoder mirrors this architecture symmetrically.
The model is trained by maximizing the Evidence Lower BOund (ELBO):
\begin{equation}
  \mathcal{L}_{\mathrm{ELBO}}
    = \mathbb{E}_{q_\phi}\bigl[\log p_\theta(\mathbf{x} \mid \mathbf{z})\bigr]
      - \beta_{\mathrm{KL}}\,
        D_{\mathrm{KL}}\!\bigl(q_\phi(\mathbf{z} \mid \mathbf{x})
                               \;\|\; p(\mathbf{z})\bigr),
  \label{eq:elbo}
\end{equation}
with $\beta_{\mathrm{KL}} = 1.0$ (standard VAE) and isotropic Gaussian prior
$p(\mathbf{z}) = \mathcal{N}(\mathbf{0}, \mathbf{I})$.
Training proceeds for 50 epochs with batch size 512 and Adam
optimizer~\cite{Kingma2014} (learning rate $0.001$).
Early stopping on the validation ELBO yields the best checkpoint at validation
loss $0.2216$.
The per-sample reconstruction error
$r(\mathbf{x}) = \|\mathbf{x} - \hat{\mathbf{x}}\|_2$
serves as the primary anomaly score in Module~A.

\paragraph{HDBSCAN density clustering.}
The VAE mean embedding $\boldsymbol{\mu}_\phi(\mathbf{x})$ is passed to
Hierarchical DBSCAN (HDBSCAN)~\cite{Campello2013} with
\texttt{min\_cluster\_size}$= 100$ and \texttt{min\_samples}$= 10$
under Euclidean metric.
HDBSCAN is preferred over $k$-means or Gaussian Mixture Models for three
reasons: (i) it discovers clusters of arbitrary shape in the latent manifold;
(ii) it explicitly models noise points rather than forcing them into the
nearest cluster; and (iii) its soft cluster membership probabilities are
available for downstream weighting.
The procedure yields \textbf{135 clusters} and
\num{63390} noise points (\SI{21.1}{\percent} of the corpus), with
normalized mutual information $\mathrm{NMI} = 0.112$ against the ground-truth
transaction-type labels—a meaningful signal given that clusters are learned
purely from graph topology and feature distributions, not label supervision.

\paragraph{Wasserstein temporal drift monitoring.}
To detect distributional shift in the transaction stream we partition
transactions into 30 consecutive monthly windows and compute the
Wasserstein-$1$ distance~\cite{Villani2009} between the cluster-label
distributions of adjacent windows:
\begin{equation}
  W_1(\mu_t, \mu_{t+1})
    = \inf_{\gamma \in \Gamma(\mu_t, \mu_{t+1})}
        \int \|c - c'\| \, d\gamma(c, c'),
  \label{eq:w1}
\end{equation}
where $\mu_t$ is the empirical cluster-assignment distribution in window $t$.
A drift alert is raised whenever $W_1(\mu_t, \mu_{t+1}) > \delta = 0.1$.
Across the 29 adjacent window pairs, all 29 pairs trigger the alert, with
$W_1$ scores ranging from $0.179$ to $50.745$, indicating substantial and
sustained distributional shift in Bitcoin transaction patterns over the
study period.
This persistent drift motivates SPECTRA's modular architecture: Module~C can
be periodically retrained in isolation without disrupting the other modules.

% ---------------------------------------------------------------------------
\subsection{Module P — Bayesian Trust Scoring}
\label{sec:module_p}

Module~P maintains a per-address trust score that is updated online as new
transactions are observed, providing a continuous and interpretable
credibility metric for each blockchain participant.

\paragraph{Beta-Bernoulli model.}
For each address $v$ we maintain a Beta posterior over the latent trust
parameter $\tau_v \in [0,1]$:
\begin{equation}
  p(\tau_v \mid \mathcal{D}_v)
    = \mathrm{Beta}\!\left(\alpha_v,\, \beta_v\right),
  \label{eq:beta_posterior}
\end{equation}
initialized with the weakly informative prior
$\mathrm{Beta}(2, 2)$ (mean $= 0.5$, mode at $0.5$, unimodal).
This prior encodes epistemic humility: absent any observed transactions, all
addresses are considered equally trustworthy.

\paragraph{Online posterior update.}
Upon observing transaction $\mathbf{x}$ at address $v$, the likelihood
of \emph{trustworthy behavior} is defined as
\begin{equation}
  \ell(\mathbf{x}) = \exp\!\bigl(-r(\mathbf{x})\bigr),
  \quad r(\mathbf{x}) = \|\mathbf{x} - \hat{\mathbf{x}}\|_2,
  \label{eq:likelihood}
\end{equation}
where $r(\mathbf{x})$ is the VAE reconstruction error from Module~C.
Low reconstruction error implies that the transaction conforms to the learned
distribution of typical activity for addresses in its cluster—high fidelity
behavior earns high trust.
The Beta parameters are updated via conjugate closed-form update:
\begin{equation}
  \alpha_v \leftarrow \alpha_v + \ell(\mathbf{x}), \quad
  \beta_v  \leftarrow \beta_v  + 1 - \ell(\mathbf{x}).
  \label{eq:beta_update}
\end{equation}
The posterior mean $\hat{\tau}_v = \alpha_v / (\alpha_v + \beta_v)$
is reported as the trust certificate $\tau_v$.

\paragraph{Population statistics.}
Across the \num{30000} hub addresses, the population trust distribution has
mean $\mu_\tau = 0.634$ and standard deviation $\sigma_\tau = 0.105$.
The right-skewed distribution—with a long low-trust tail—reflects the
Bitcoin ecosystem: a majority of addresses exhibit predictable, repetitive
behavioral patterns (exchange hot wallets reusing identical fee structures)
while a minority show highly variable, high-entropy activity consistent
with mixing or layering behavior.

% ---------------------------------------------------------------------------
\subsection{Module T — Markov Chain Behavioral Profiling}
\label{sec:module_t}

Module~T models the cluster-state trajectory of each address as a discrete
Markov chain, capturing behavioral patterns that unfold over sequences of
transactions rather than in individual events.

\paragraph{State space and transition estimation.}
The state space $\mathcal{S}$ consists of the $C = 135$ HDBSCAN cluster
labels plus one absorbing noise state, giving $|\mathcal{S}| = 136$.
From $\num{55678}$ multi-transaction address sequences we estimate the
maximum-likelihood transition matrix
$\mathbf{P} \in \mathbb{R}^{136 \times 136}$, where
\begin{equation}
  P_{ij} = \frac{n_{ij} + \varepsilon}{\sum_{j'} (n_{ij'} + \varepsilon)},
  \quad \varepsilon = 10^{-6},
  \label{eq:markov_transition}
\end{equation}
with $n_{ij}$ the observed count of transitions from state $i$ to state $j$
and $\varepsilon = 10^{-6}$ Laplace smoothing to avoid zero-probability
transitions for unobserved pairs.
The stationary distribution $\boldsymbol{\pi} = [\pi_1, \ldots, \pi_{136}]$
is computed via power iteration with tolerance $10^{-8}$.

\paragraph{Transition anomaly scoring.}
The anomaly score for a single transition $s_{t-1} \to s_t$ is the negative
log transition probability:
\begin{equation}
  \mathrm{score}(s_{t-1}, s_t) = -\log P_{s_{t-1}, s_t}.
  \label{eq:markov_score}
\end{equation}
High scores flag rare—and therefore anomalous—behavioral transitions, such as
an address that consistently transacts within one cluster (e.g., exchange
deposit cluster) suddenly appearing in the mixing-service cluster.
Across all observed transitions the population mean anomaly score is
$\bar{m} = 0.875$; Module~A flags an address if its per-transition score
exceeds $\bar{m} + 2\sigma_m$.

% ---------------------------------------------------------------------------
\subsection{Module R — GraphSAGE Transaction Recognition}
\label{sec:module_r}

Module~R performs multi-class transaction-type classification using a
two-layer GraphSAGE network~\cite{Hamilton2017} that aggregates
neighborhood information from the Bitcoin address graph.

\paragraph{Input feature concatenation.}
The input feature vector for address node $v$ is assembled as
\begin{equation}
  \mathbf{f}_v =
    \bigl[\,
      \mathbf{u}_v \;\|\;
      \boldsymbol{\mu}_v \;\|\;
      \tau_v \;\|\;
      m_v \;\|\;
      \mathbf{e}_v
    \,\bigr]
    \in \mathbb{R}^{d_f},
  \label{eq:feat_concat}
\end{equation}
where $\mathbf{u}_v \in \mathbb{R}^{16}$ is the spectral-behavioral embedding
from Module~S, $\boldsymbol{\mu}_v \in \mathbb{R}^{32}$ is the VAE latent
mean from Module~C, $\tau_v \in \mathbb{R}$ is the Bayesian trust score from
Module~P, $m_v \in \mathbb{R}$ is the Markov anomaly score from Module~T, and
$\mathbf{e}_v \in \mathbb{R}^{136}$ is the one-hot HDBSCAN cluster assignment
(including the noise label).
The total input dimension is $d_f = 16 + 32 + 1 + 1 + 136 = 186$.

\paragraph{GraphSAGE architecture.}
GraphSAGE~\cite{Hamilton2017} learns a neighborhood aggregation
function that generalizes inductively to unseen nodes:
\begin{equation}
  \mathbf{h}_v^{(l)}
    = \sigma\!\left(
        \mathbf{W}^{(l)}
        \cdot
        \mathrm{CONCAT}\!\left(
          \mathbf{h}_v^{(l-1)},\;
          \mathrm{MEAN}_{u \in \mathcal{N}(v)}\!\left\{
            \mathbf{h}_u^{(l-1)}
          \right\}
        \right)
      \right),
  \label{eq:graphsage}
\end{equation}
where $\sigma$ is the ReLU activation and $\mathcal{N}(v)$ is a fixed-size
sampled neighborhood.
We use two GraphSAGE layers with hidden dimension 128 and output embedding
dimension 64, followed by a linear classification head over six transaction
types.
Dropout with rate $0.3$ is applied after each layer.

\paragraph{Rationale for GraphSAGE over GAT.}
Graph Attention Networks (GAT)~\cite{Velickovic2018} learn attention
weights over neighbors but impose quadratic attention-head complexity per
node.
For the sparse, non-attributed Bitcoin address graph three considerations
favor GraphSAGE: (i) \emph{scalability} — neighborhood sampling in GraphSAGE
scales sub-linearly with graph size, critical for the \num{30000}-node pruned
subgraph and its highly skewed degree distribution; (ii) \emph{absence of
node attributes} — attention mechanisms derive most benefit from rich semantic
attributes (e.g., text or image features); on structural Bitcoin graphs the
gain is marginal while training cost increases significantly; and (iii)
\emph{inductive generalization} — GraphSAGE's mean aggregator generalizes to
unseen addresses appearing in the live transaction stream without retraining,
a critical operational requirement for real-time monitoring.

\paragraph{Training protocol.}
The model is trained for a maximum of 100 epochs using Adam with learning
rate $0.001$ and weight decay $10^{-4}$, with early stopping
(patience $= 15$ epochs) on the validation AUC-ROC.
Training terminates at epoch 18, yielding a total train time of
\SI{2.09}{\second} on the pruned subgraph—confirming that the modular
spectral preprocessing in Module~S has already extracted the dominant
structural signal, leaving a compact learning problem for the GNN.

% ---------------------------------------------------------------------------
\subsection{Module A — Anomaly and Drift Integration}
\label{sec:module_a}

Module~A is the decision layer of \textsc{SPECTRA}: it fuses the outputs of
Modules~C, P, and~T into per-address anomaly flags and drift alerts,
and emits a final trust certificate.

\paragraph{Unified anomaly flag.}
An address $v$ is flagged anomalous if \emph{any} of three complementary
conditions holds:
\begin{equation}
  \mathrm{anomaly}(v) = 1 \iff
    r(\mathbf{x}_v) > \tau_{\mathrm{anom}}
    \;\lor\;
    m_v > \bar{m} + 2\sigma_m
    \;\lor\;
    \mathcal{W}(t) > \delta,
  \label{eq:anomaly_flag}
\end{equation}
where $r(\mathbf{x}_v)$ is the VAE reconstruction error,
$\tau_{\mathrm{anom}} = 0.1694$ is the $95^{\mathrm{th}}$ percentile of the
population reconstruction error (learned from the training set),
$m_v$ is the Markov anomaly score from \eqref{eq:markov_score},
$\bar{m}$ and $\sigma_m$ are its population mean and standard deviation, and
$\mathcal{W}(t) = W_1(\mu_{t-1}, \mu_t)$ is the Wasserstein distance for
time window $t$ with threshold $\delta = 0.1$.
The disjunctive combination ensures that an address is flagged by any
evidence type; this union-of-detectors design minimizes false negatives at the
cost of a controlled increase in false positives, which is the preferred
operating point for financial forensics~\cite{Bahnsen2016}.

\paragraph{Trust certificate.}
The final trust certificate for address $v$ is the posterior mean trust score
$\hat{\tau}_v = \alpha_v / (\alpha_v + \beta_v) \in [0,1]$ from Module~P,
optionally attenuated by the anomaly flag:
$\tilde{\tau}_v = \hat{\tau}_v \cdot (1 - \mathrm{anomaly}(v))$.
This scalar certificate is the primary output delivered to downstream
compliance or forensic systems.

\paragraph{Drift alert.}
In addition to per-address flags, Module~A emits a binary per-window drift
alert $\mathcal{A}_t = \mathbf{1}[\mathcal{W}(t) > \delta]$.
Sustained drift alerts over consecutive windows (all 29 of the 29 observable
pairs in the study period) indicate concept drift requiring retraining of
Module~C's VAE and HDBSCAN components.
Because the Markov transition matrix (Module~T) and GraphSAGE weights
(Module~R) are conditioned on the cluster state space defined by Module~C,
a cluster re-labeling event requires propagating updated cluster assignments
through Modules~P, T, and~R; the modular architecture ensures this can be
done without restarting the pipeline.

% ---------------------------------------------------------------------------
\subsection{Unified Algorithm}
\label{sec:algorithm}

Algorithm~\ref{alg:spectra} presents the complete \textsc{SPECTRA} pipeline
in pseudocode, showing the sequential data flow across all seven modules.
Notation follows~\eqref{eq:lsym}--\eqref{eq:anomaly_flag}.

% ---------------------------------------------------------------------------
\begin{algorithm}[t]
\caption{\textsc{SPECTRA} — Spectral-Probabilistic Ensemble Clustering for
         Transaction Recognition and Authentication}
\label{alg:spectra}
\begin{algorithmic}[1]
\Require Raw transaction dataset $\mathcal{D}$ (\num{300000} transactions);
         hyperparameters $k, r, \beta_{\mathrm{KL}}, \delta, \varepsilon$
\Ensure  Per-address trust certificate $\tilde{\tau}_v$,
         anomaly flag $\mathrm{anomaly}(v)$,
         transaction-type prediction $\hat{y}_v$,
         drift alert $\mathcal{A}_t$

\Statex
\Statex \textbf{Module S — Spectral Graph Construction}
\State Build directed weighted address graph
       $\mathcal{G}=(\mathcal{V},\mathcal{E},\mathbf{W})$
       from $\mathcal{D}$\quad
       ($|\mathcal{V}|=\num{2870000}$, $|\mathcal{E}|=\num{28600000}$)
\State Prune $\mathcal{G}$ to top-\num{30000} nodes by weighted in-degree
       $\Rightarrow \mathcal{G}'$
\State Compute symmetric normalized Laplacian
       $\mathbf{L}_{\mathrm{sym}} = \mathbf{I} -
        \mathbf{D}^{-1/2}\mathbf{A}\mathbf{D}^{-1/2}$
       on $\mathcal{G}'$
\State Extract $k=10$ spectral embeddings
       $\mathbf{U} \in \mathbb{R}^{N \times 10}$
       via randomized SVD (\texttt{n\_iter}$=10$, \texttt{n\_oversamples}$=20$)
\State Compute 6 behavioral features per node; concatenate with $\mathbf{U}$
       $\Rightarrow \mathbf{X}_{\mathrm{node}} \in \mathbb{R}^{N \times 16}$

\Statex
\Statex \textbf{Module E — Information-Theoretic Feature Ranking}
\State Compute Shannon entropy $H(X_j)$ for each raw feature $X_j$;
       discard low-entropy features
\State Estimate $I(X_j; Y)$ via $k$-NN mutual information;
       retain top-20 features
       $\Rightarrow \mathbf{x}_{\mathrm{tx}} \in \mathbb{R}^{M \times 20}$
\State Fit NMF ($r=10$, $\alpha_W=0.1$, $\ell_1$-ratio$=0.5$)
       $\Rightarrow \mathbf{H} \in \mathbb{R}^{M \times 10}$;
       concatenate $[\mathbf{x}_{\mathrm{tx}} \| \mathbf{H}]$

\Statex
\Statex \textbf{Module C — VAE + HDBSCAN Clustering}
\State Train VAE (hidden$=[128,64]$, latent$=32$, $\beta_{\mathrm{KL}}=1.0$,
       dropout$=0.2$, 50 epochs, lr$=0.001$) on $\mathbf{x}_{\mathrm{tx}}$
\State Encode all transactions:
       $(\boldsymbol{\mu}_v, \boldsymbol{\sigma}^2_v)
        \leftarrow q_\phi(\mathbf{z} \mid \mathbf{x}_v)$;
       record reconstruction error $r(\mathbf{x}_v)$
\State Cluster $\{\boldsymbol{\mu}_v\}$ with HDBSCAN
       (\texttt{min\_cluster\_size}$=100$, \texttt{min\_samples}$=10$)
       $\Rightarrow$ 135 clusters + noise; labels $\{c_v\}$
\For{each adjacent monthly window pair $(t-1, t)$, $t=1,\ldots,29$}
    \State Compute $\mathcal{W}(t) = W_1(\mu_{t-1}, \mu_t)$;
           set $\mathcal{A}_t = \mathbf{1}[\mathcal{W}(t) > \delta]$
\EndFor

\Statex
\Statex \textbf{Module P — Bayesian Trust Scoring}
\State Initialize $(\alpha_v, \beta_v) \leftarrow (2, 2)$ for all $v$
\For{each transaction $\mathbf{x}$ at address $v$ (in temporal order)}
    \State $\ell \leftarrow \exp(-r(\mathbf{x}))$
    \State $\alpha_v \leftarrow \alpha_v + \ell$;\quad
           $\beta_v \leftarrow \beta_v + (1-\ell)$
\EndFor
\State $\hat{\tau}_v \leftarrow \alpha_v / (\alpha_v + \beta_v)$
       for all $v$

\Statex
\Statex \textbf{Module T — Markov Chain Behavioral Profiling}
\State Build transition count matrix $\mathbf{N} \in \mathbb{R}^{136 \times 136}$
       from address cluster-state sequences
\State Estimate $\mathbf{P}$:
       $P_{ij} \leftarrow (n_{ij} + \varepsilon) /
                           (\sum_{j'} n_{ij'} + 136\varepsilon)$,\;
       $\varepsilon = 10^{-6}$
\State Compute stationary distribution $\boldsymbol{\pi}$ via power iteration
       (tol $=10^{-8}$)
\State $m_v \leftarrow -\log P_{c_{v,t-1},\, c_{v,t}}$
       for each transition of address $v$

\Statex
\Statex \textbf{Module R — GraphSAGE Transaction Recognition}
\State Assemble node feature matrix:
       $\mathbf{f}_v =
        [\mathbf{X}_{\mathrm{node},v} \;\|\;
         \boldsymbol{\mu}_v \;\|\;
         \hat{\tau}_v \;\|\;
         m_v \;\|\;
         \mathbf{e}_{c_v}]
        \in \mathbb{R}^{186}$
\State Train two-layer GraphSAGE (hidden$=128$, embed$=64$, dropout$=0.3$,
       lr$=0.001$, wd$=10^{-4}$, patience$=15$) on $\mathcal{G}'$
       with node features $\mathbf{F}$ and transaction-type labels
\State Infer $\hat{y}_v \leftarrow \mathrm{argmax}\;\mathrm{softmax}
       (\mathbf{h}_v^{(2)}\mathbf{W}_{\mathrm{cls}})$ for all $v$

\Statex
\Statex \textbf{Module A — Anomaly and Drift Integration}
\State Compute population statistics:
       $\tau_{\mathrm{anom}} \leftarrow \mathrm{perc}_{95}(\{r(\mathbf{x}_v)\})$;\quad
       $\bar{m}, \sigma_m \leftarrow \mathrm{mean}, \mathrm{std}(\{m_v\})$
\For{each address $v$}
    \State $\mathrm{anomaly}(v) \leftarrow
            \mathbf{1}\!\left[
              r(\mathbf{x}_v) > \tau_{\mathrm{anom}}
              \;\lor\;
              m_v > \bar{m} + 2\sigma_m
              \;\lor\;
              \mathcal{A}_{t(v)} = 1
            \right]$
    \State $\tilde{\tau}_v \leftarrow \hat{\tau}_v \cdot
            (1 - \mathrm{anomaly}(v))$
\EndFor
\State \Return $\{\hat{y}_v\}$, $\{\tilde{\tau}_v\}$,
               $\{\mathrm{anomaly}(v)\}$, $\{\mathcal{A}_t\}$
\end{algorithmic}
\end{algorithm}
% ---------------------------------------------------------------------------

\subsection{Complexity Analysis}
\label{sec:complexity}

Table~\ref{tab:complexity} summarizes the computational complexity of each
module.
The dominant cost is Module~S: randomized SVD on $\mathbf{L}_{\mathrm{sym}}$
runs in $\mathcal{O}(N k \log k)$ time~\cite{Halko2011}, far cheaper
than the $\mathcal{O}(N^2)$ dense eigensolver.
Module~C's VAE training is $\mathcal{O}(M \cdot B^{-1} \cdot E_{\max})$
where $B = 512$ is the batch size and $E_{\max} = 50$ the epoch budget;
HDBSCAN is $\mathcal{O}(M \log M)$ for the mutual reachability graph
construction phase.
Module~T's transition matrix estimation is $\mathcal{O}(|\mathcal{S}|^2)$
for matrix storage and $\mathcal{O}(|\mathcal{S}|^2 / (1-\lambda_2))$ for
power iteration convergence, where $\lambda_2$ is the second eigenvalue of
$\mathbf{P}$.
Module~R's GraphSAGE forward pass is $\mathcal{O}(N \cdot d_f \cdot L \cdot S)$
where $L = 2$ is the number of layers and $S = 25$ the neighborhood sample
size; early stopping at epoch 18 further limits wall-clock training time to
\SI{2.09}{\second}.
End-to-end inference for a new address (given frozen modules) is dominated by
the VAE encoder ($\mathcal{O}(d_f \cdot 128)$) and GraphSAGE forward pass
($\mathcal{O}(d_f \cdot S \cdot 128)$), both operating in sub-millisecond
wall-clock time on a single CPU core.

\begin{table}[t]
  \caption{Per-module complexity and key hyperparameters of \textsc{SPECTRA}.
           $N=\num{30000}$ (pruned nodes), $M=\num{300000}$ (transactions),
           $k=10$ (spectral dims), $C=136$ (cluster states),
           $d_f=186$ (GNN input dim).}
  \label{tab:module_summary}
  \centering
  \renewcommand{\arraystretch}{1.15}
  \begin{tabular}{@{}lllp{3.6cm}@{}}
    \toprule
    \textbf{Module} & \textbf{Output} & \textbf{Complexity} & \textbf{Key hyperparameters} \\
    \midrule
    S — Spectral  & $\mathbf{X}_{\mathrm{node}} \!\in\! \mathbb{R}^{N\times 16}$
                  & $\mathcal{O}(Nk\log k)$
                  & $k{=}10$, \texttt{n\_iter}$=10$ \\
    E — Entropy   & $\mathbf{x}_{\mathrm{tx}} \!\in\! \mathbb{R}^{M\times 20}$
                  & $\mathcal{O}(MF)$
                  & top-20 MI, NMF $r{=}10$ \\
    C — Cluster   & $\{c_v\}$, $\{\boldsymbol{\mu}_v\}$, $\{\mathcal{A}_t\}$
                  & $\mathcal{O}(M\log M)$
                  & latent$=32$, MCS$=100$ \\
    P — Probabilistic & $\hat{\tau}_v \!\in\! [0,1]$
                  & $\mathcal{O}(M)$
                  & $\mathrm{Beta}(2,2)$ prior \\
    T — Markov    & $m_v \!\in\! \mathbb{R}_{+}$
                  & $\mathcal{O}(C^2)$
                  & $|\mathcal{S}|{=}136$, $\varepsilon{=}10^{-6}$ \\
    R — Recognition & $\hat{y}_v \!\in\! \{1,\ldots,6\}$
                  & $\mathcal{O}(NLSd_f)$
                  & hidden$=128$, dropout$=0.3$ \\
    A — Anomaly   & $\mathrm{anomaly}(v)$, $\tilde{\tau}_v$
                  & $\mathcal{O}(N)$
                  & $\tau_{\mathrm{anom}}{=}0.1694$, $\delta{=}0.1$ \\
    \bottomrule
  \end{tabular}
\end{table}

\begin{table}[t]
  \caption{Top-20 transaction features ranked by mutual information
           $I(X_j; Y)$ with transaction-type label $Y$.
           All values rounded to three decimal places.}
  \label{tab:mi_top20}
  \centering
  \renewcommand{\arraystretch}{1.12}
  \begin{tabular}{@{}clc@{}}
    \toprule
    \textbf{Rank} & \textbf{Feature} & $I(X_j;\,Y)$ \\
    \midrule
    1  & \texttt{weight}                   & 1.320 \\
    2  & \texttt{vsize}                    & 1.204 \\
    3  & \texttt{size}                     & 1.193 \\
    4  & \texttt{output\_count}            & 1.171 \\
    5  & \texttt{input\_count}             & 0.726 \\
    6  & \texttt{fee\_rate\_sat\_per\_byte}& 0.494 \\
    7  & \texttt{fee}                      & 0.471 \\
    8  & \texttt{value\_difference}        & 0.470 \\
    9  & \texttt{fee\_rate\_sat\_per\_vbyte}& 0.410 \\
    10 & \texttt{total\_output\_value}     & 0.380 \\
    \multicolumn{3}{c}{\textit{(ranks 11--20 reported in supplementary material)}} \\
    \bottomrule
  \end{tabular}
\end{table}

% =============================================================================
% SECTION V — EXPERIMENTAL SETUP
% =============================================================================

% =============================================================================
%  SPECTRA — Section V: Experimental Setup
%  Target: IEEE Transactions on Information Forensics & Security
%  Author: Sagar Korde
% =============================================================================

\section{Experimental Evaluation}
\label{sec:experiments}

% ─── V-A  Dataset ─────────────────────────────────────────────────────────────

\subsection{Dataset}
\label{subsec:dataset}

We evaluate SPECTRA on the CryptoGraphNet dataset~\cite{cryptographnet2024},
a publicly available corpus of Bitcoin transactions extracted from blocks
mined between January 2009 and December 2023.
The full dataset contains \num{5884387} transactions with \num{53} fields per
record, including raw structural attributes (input/output counts, values, fees,
script types), derived behavioral indicators, and assigned transaction-type labels.

We draw a stratified sample of $N = \num{300000}$ transactions distributed
uniformly across six transaction types:
Standard ($n = \num{50000}$),
Peer-to-Peer ($n = \num{50000}$),
Consolidation ($n = \num{50000}$),
Distribution ($n = \num{50000}$),
Batch-Payment ($n = \num{50000}$), and
CoinJoin-like ($n = \num{50000}$).
The stratified sampling ensures balanced class representation despite the
long-tailed distribution in the full corpus (Standard transactions dominate).

The same \num{300000} transactions serve triple duty:
(i)~as the ML classification dataset (70/15/15 stratified train/val/test split),
(ii)~as the source graph for spectral embedding ($N_G = 300\,000$,
pruned to $|\mathcal{V}| = 30\,000$ hub nodes by undirected degree), and
(iii)~as the observation window for Wasserstein drift detection and
Markov chain fitting.
Using the same transaction pool for both graph construction and classification
is a deliberate design choice that maximizes address coverage: the top-30\,000
hub addresses selected by degree represent the most-connected actors
(exchanges, mining pools, mixers) and appear in the highest fraction of
all transactions.

% ─── V-B  Feature Engineering ────────────────────────────────────────────────

\subsection{Feature Engineering}
\label{subsec:features}

We begin with a raw feature matrix $\mathbf{X}_{\text{raw}} \in \mathbb{R}^{N \times 31}$
obtained by dropping identifier columns (\texttt{txid}, address arrays,
script type strings) and always-zero Boolean flags
(\texttt{is\_self\_transfer}, \texttt{has\_taproot}, etc.)\ that carry no
discriminative information.
From the remaining 31 numeric features we compute Shannon entropy per column
and select the top $k = 20$ features by Mutual Information (MI) with the
transaction-type label~\cite{Kraskov2004}.

The final \emph{temporally-fair} feature set $\mathbf{x}^{\dagger}$
is obtained by additionally removing the three absolute temporal identifiers
\texttt{block\_height}, \texttt{block\_time}, and \texttt{year}
(see Section~\ref{subsec:label_provenance}).
The 20 features selected after this temporal filtering are, in descending
MI order: \texttt{weight}, \texttt{vsize}, \texttt{size}, \texttt{output\_count},
\texttt{input\_count}, \texttt{fee\_rate\_sat\_per\_byte}, \texttt{fee},
\texttt{value\_difference}, \texttt{fee\_rate\_sat\_per\_vbyte},
\texttt{total\_output\_value}, \texttt{avg\_output\_value},
\texttt{total\_input\_value}, \texttt{avg\_input\_value},
\texttt{input\_output\_ratio}, \texttt{has\_op\_return}, \texttt{rbf\_enabled},
\texttt{version}, \texttt{week\_of\_year}, \texttt{sample\_size},
and \texttt{locktime}.
All features are standardized to zero mean and unit variance on the training
set; statistics are held fixed for validation and test sets.

Additionally, Non-Negative Matrix Factorization (NMF) with $r = 10$ components
decomposes the feature matrix to extract latent behavioral patterns
(Module~E)~\cite{Lee1999}.
The NMF basis vectors $\mathbf{H} \in \mathbb{R}^{10 \times 20}$ are reported
as behavioral profiles in Fig.~\ref{fig:nmf_patterns}.

% ─── V-C  Graph Construction ─────────────────────────────────────────────────

\subsection{Graph Construction and Spectral Embedding}
\label{subsec:graph}

We construct a directed weighted graph
$\mathcal{G} = (\mathcal{V}, \mathcal{E}, \mathbf{W})$
where $\mathcal{V}$ is the set of unique wallet addresses,
$(u, v) \in \mathcal{E}$ iff address $u$ sent BTC to address $v$ in at
least one transaction, and edge weight $w_{uv}$ accumulates the total
transferred value (BTC).
Bitcoin addresses are parsed from VARCHAR[] fields using semicolon
de-serialization as per the CryptoGraphNet data format.

From the full \num{300000}-transaction graph (approximately $3 \times 10^6$
raw nodes), we retain the top $|\mathcal{V}| = 30\,000$ nodes by
undirected degree.
These high-degree hubs correspond to exchanges, mining pools, and mixing
services, which carry the most structural signal for the normalized
Laplacian eigenvectors.
The resulting pruned graph contains approximately $5 \times 10^6$ directed
edges, representing the dense interconnection among the most active
blockchain actors.

The normalized graph Laplacian
$\mathbf{L}_{\text{sym}} = \mathbf{D}^{-1/2}(\mathbf{D} - \mathbf{A}_{\text{sym}})\mathbf{D}^{-1/2}$,
where $\mathbf{A}_{\text{sym}} = (\mathbf{A} + \mathbf{A}^\top)/2$
symmetrizes the directed adjacency matrix, is decomposed via randomized SVD
(sklearn, $n\_\text{iter}=10$, $n\_\text{oversamples}=20$)~\cite{Halko2011}
to obtain the bottom $k = 10$ non-trivial eigenvectors
$\boldsymbol{\Phi} \in \mathbb{R}^{30000 \times 10}$.
Randomized SVD replaces the ARPACK eigensolver used in prior
work~\cite{Kipf2017}, resolving convergence failures on this graph's
density and Python~3.12 GIL-related threading crashes in ARPACK's
FORTRAN backend.

Each node's spectral features are augmented with six behavioral statistics
computed from its transaction history:
fan-in ratio, transaction velocity (tx/day), average received BTC,
average sent BTC, mean and standard deviation of inter-transaction gap.
The resulting node feature matrix is
$\mathbf{X}_{\text{node}} \in \mathbb{R}^{30000 \times 16}$.

% ─── V-D  Baseline Implementations ──────────────────────────────────────────

\subsection{Baseline Implementations}
\label{subsec:baselines}

We implement ten IEEE-mandatory baselines under identical train/val/test splits
and standardized feature inputs.
All methods operate on either $\mathbf{X}^{\dagger}_{\text{scaled}} \in
\mathbb{R}^{N \times 20}$ (tabular track) or the graph topology
$(\mathcal{G}, \boldsymbol{\Phi})$ (graph-topology track).

\begin{enumerate}[leftmargin=*, label=\textbf{B\arabic*}]

  \item \textbf{K-Means}~\cite{Lloyd1982}: $k = 6$ clusters, $k$-means++
    initialization. Evaluated by Silhouette coefficient, Davies-Bouldin index,
    Calinski-Harabasz score, and Normalized Mutual Information (NMI).

  \item \textbf{DBSCAN}~\cite{Ester1996}: $\varepsilon = 0.5$,
    $\textit{minPts} = 10$. Noise points excluded from cluster metrics.

  \item \textbf{Isolation Forest}~\cite{Liu2008}: contamination $= 0.05$.
    Anomaly scores binarized at the 95th-percentile threshold.

  \item \textbf{Autoencoder + K-Means}: Two-layer MLP autoencoder
    (hidden dims $[64, 32]$, 30 epochs) followed by $k$-means on the
    32-dimensional bottleneck. Represents deep unsupervised baseline without
    probabilistic guarantees.

  \item \textbf{Plain HDBSCAN}~\cite{McInnes2017}: $\textit{min\_cluster\_size}
    = 100$, $\textit{min\_samples} = 10$, applied directly to
    $\mathbf{X}^{\dagger}_{\text{scaled}}$ (no VAE encoding).
    Isolates the contribution of VAE latent-space projection.

  \item \textbf{XGBoost}~\cite{Chen2016}: $300$ trees, max depth $6$,
    learning rate $0.05$. The strongest non-DL tabular baseline.

  \item \textbf{EvolveGCN}~\cite{Pareja2020}: Simplified EvolveGCN-O variant
    with GCN layers and LSTMCell weight evolution.
    Applied in tabular mode (no edge index at inference) as in
    Jha~\emph{et al.}~\cite{Jha2023}.

  \item \textbf{BERT4ETH-BTC}~\cite{He2023}: Transformer encoder adapted
    to Bitcoin address sequences.  Projection layer maps raw features to
    embedding space; CLS-token representation fed to softmax head.

  \item \textbf{Graph Autoencoder (GAE)}~\cite{Kipf2016b}: GCN encoder
    with inner-product decoder, anomaly score = reconstruction error.
    Bitcoin application per Liu~\emph{et al.}~\cite{Liu2024}.

  \item \textbf{AA-GNN}~\cite{Fan2024}: GAT encoder with anomaly-aware
    residual head trained via Focal Loss~\cite{Lin2017}.

\end{enumerate}

% ─── V-E  Evaluation Metrics ──────────────────────────────────────────────────

\subsection{Evaluation Metrics}
\label{subsec:metrics}

For \textbf{classification} baselines (B6--B10, SPECTRA): accuracy,
macro-F$_1$, weighted-F$_1$, macro-precision, macro-recall,
AUC-ROC (one-vs-rest, macro-averaged), AUC-PR (macro-averaged),
and Matthews Correlation Coefficient (MCC).

For \textbf{clustering} baselines (B1--B5, B9):
Silhouette coefficient ($\uparrow$), Davies-Bouldin index ($\downarrow$),
Calinski-Harabasz score ($\uparrow$), and NMI ($\uparrow$).

For \textbf{SPECTRA}'s additional outputs: Wasserstein-$1$ drift score
per time window, per-address Bayesian trust score $\tau \in [0,1]$,
and Markov chain stationarity convergence.

All experiments use fixed seed $= 42$.
Hardware: NVIDIA RTX 4060 Laptop GPU (8 GB VRAM), Intel Core i7,
16 GB RAM, Windows 11, Python 3.10, PyTorch 2.7, PyG 2.6.

% ─── References used above ───────────────────────────────────────────────────
% Add to bibliography:
%
% \bibitem{Kraskov2004} A.~Kraskov, H.~Stögbauer, P.~Grassberger,
%   ``Estimating mutual information,'' Phys.\ Rev.\ E, vol.~69, 2004.
%
% \bibitem{Lee1999} D.~D.~Lee and H.~S.~Seung,
%   ``Learning the parts of objects by non-negative matrix factorization,''
%   Nature, vol.~401, pp.~788--791, 1999.
%
% \bibitem{Halko2011} N.~Halko, P.-G.~Martinsson, J.~A.~Tropp,
%   ``Finding structure with randomness: Probabilistic algorithms for
%   constructing approximate matrix decompositions,''
%   SIAM Rev., vol.~53, no.~2, pp.~217--288, 2011.
%
% \bibitem{Lloyd1982} S.~Lloyd, ``Least squares quantization in PCM,''
%   IEEE Trans.\ Inf.\ Theory, vol.~28, pp.~129--137, 1982.
%
% \bibitem{Ester1996} M.~Ester et al., ``A density-based algorithm for
%   discovering clusters in large spatial databases,'' KDD, 1996.
%
% \bibitem{Liu2008} F.~T.~Liu, K.~M.~Ting, Z.-H.~Zhou,
%   ``Isolation Forest,'' ICDM, 2008.
%
% \bibitem{McInnes2017} L.~McInnes, J.~Healy, S.~Astels,
%   ``hdbscan: Hierarchical density based clustering,''
%   JOSS, vol.~2, no.~11, 2017.
%
% \bibitem{Chen2016} T.~Chen and C.~Guestrin,
%   ``XGBoost: A scalable tree boosting system,'' KDD, 2016.
%
% \bibitem{Pareja2020} A.~Pareja et al.,
%   ``EvolveGCN: Evolving graph convolutional networks for dynamic graphs,''
%   AAAI, 2020.
%
% \bibitem{Jha2023} S.~Jha et al.,
%   ``Temporal graph network for Bitcoin transaction analysis,''
%   IEEE Trans.\ Inf.\ Forensics Security, vol.~18, 2023.
%
% \bibitem{He2023} Z.~He et al.,
%   ``BERT4ETH: A pre-trained transformer for Ethereum transaction behavior
%   sequence modeling,'' WWW, 2023.
%
% \bibitem{Kipf2016b} T.~N.~Kipf and M.~Welling,
%   ``Variational graph auto-encoders,'' NIPS Workshop, 2016.
%
% \bibitem{Liu2024} Y.~Liu et al.,
%   ``Graph autoencoder for blockchain anomaly detection,''
%   IEEE Trans.\ Dependable Secure Comput., 2024.
%
% \bibitem{Fan2024} S.~Fan et al.,
%   ``Anomaly-aware graph neural networks for blockchain transaction
%   fraud detection,'' IEEE Trans.\ Ind.\ Inform., 2024.
%
% \bibitem{Lin2017} T.-Y.~Lin et al.,
%   ``Focal loss for dense object detection,'' ICCV, 2017.

% ----- V-B: Label Provenance (sub-section, insert after V-A) -----------------

% =============================================================================
%  SPECTRA — Paper Section Draft
%  Label Provenance & Evaluation Methodology
%  Target: IEEE Transactions on Information Forensics & Security
%  Author: Sagar Korde
% =============================================================================
%
%  INSERT this block inside Section V (Experimental Evaluation),
%  as a subsection before Table I (Baseline Comparison).
% =============================================================================

% ─── V-B  Label Provenance and Evaluation Methodology ────────────────────────

\subsection{Label Provenance and Evaluation Methodology}
\label{subsec:label_provenance}

\subsubsection{Dataset and Transaction Type Labels}
We evaluate on the CryptoGraphNet dataset~\cite{cryptographnet2024}, which contains
\num{300000} stratified Bitcoin transactions sampled across six mutually exclusive
transaction types: Standard, Peer-to-Peer (P2P), Consolidation, Distribution,
Batch-Payment, and CoinJoin-like.
The six class labels are assigned via deterministic heuristics applied to
structural transaction attributes---specifically, the number of distinct input
and output addresses, value concentration ratios, and fee-rate patterns.
For example, a \emph{Consolidation} transaction is defined as one with
$|\mathit{inputs}| \geq 3$ and $|\mathit{outputs}| = 1$, while a
\emph{Batch-Payment} requires $|\mathit{outputs}| \geq 5$~\cite{cryptographnet2024}.

\subsubsection{Label--Feature Entanglement}
This construction introduces a structural entanglement between labels and tabular
features: because the class membership rules are algebraic functions of
$|\mathit{inputs}|$, $|\mathit{outputs}|$, and value statistics, any model
that observes these raw tabular features can reconstruct the labels
with near-perfect accuracy without learning any generalizable behavioral
pattern.
Table~\ref{tab:fair_baselines} confirms this: XGBoost and
BERT4ETH-BTC both achieve $\text{F}_1 = 1.000$ on the original feature set.
Critically, this degeneracy \emph{persists even after temporal identifiers
are removed} (Table~\ref{tab:fair_baselines}, ``fair'' columns), because
the definitive features are transaction-structural, not temporal.

We emphasize that this is a property of the benchmark's label-generation
procedure---a known limitation of rule-based annotation common to
blockchain datasets~\cite{Jourdan2019,Akcora2022}---and not an indication
of model overfitting.
We report tabular-classifier results for completeness per IEEE TIFS
reproducibility requirements, and clearly distinguish them from the
more informative \emph{graph-topology} track discussed below.

\subsubsection{Two-Track Evaluation Protocol}
To provide a rigorous and interpretable comparison, we partition all methods
into two evaluation tracks:

\begin{enumerate}[leftmargin=*, label=\textbf{T\arabic*}]

  \item \textbf{Tabular-feature classifiers} (XGBoost, BERT4ETH-BTC).
    These methods receive the full per-transaction feature vector
    $\mathbf{x} \in \mathbb{R}^{20}$ selected by Mutual Information ranking.
    Because $\mathbf{x}$ includes $|\mathit{inputs}|$, $|\mathit{outputs}|$,
    and value statistics from which the labels are derived, their accuracy
    constitutes an \emph{upper-bound oracle} rather than a practically
    meaningful measure of fraud-detection capability.

  \item \textbf{Graph-topology methods} (SPECTRA, EvolveGCN, AA-GNN,
    clustering baselines).
    These methods operate exclusively on the address-level transaction graph
    $\mathcal{G} = (\mathcal{V}, \mathcal{E}, \mathbf{W})$ and/or the
    spectral node embedding $\boldsymbol{\Phi} \in \mathbb{R}^{|\mathcal{V}| \times k}$
    derived from the graph Laplacian.
    Raw structural features ($|\mathit{inputs}|$, $|\mathit{outputs}|$) are
    \emph{not} directly available to these methods; the graph topology
    encodes them only implicitly through node degree distributions.
    This track constitutes the \emph{meaningful} comparison and is the
    focus of SPECTRA's contribution claims.

\end{enumerate}

\subsubsection{Temporally-Fair Feature Set}
Independent of the two-track split, we additionally ablate temporal
identifier features.
In Bitcoin's ledger, \texttt{block\_height} and \texttt{block\_time}
are monotonically increasing proxies for calendar time.
Because different transaction types were prevalent at different epochs
of Bitcoin's adoption curve (e.g., CoinJoin-like transactions became
common only after 2017, and batch payments correlate with exchange activity
peaks in 2020--2021), a model with access to these fields can exploit
distributional shift rather than structural features.
We therefore define a \emph{temporally-fair} feature set
$\mathbf{x}^{\dagger}$ by removing \texttt{block\_height},
\texttt{block\_time}, and \texttt{year}, then re-selecting the top-20
features by Mutual Information on the remaining columns.
All results are reported for both the original set $\mathbf{x}$ and
the fair set $\mathbf{x}^{\dagger}$ in Table~\ref{tab:combined_comparison}.
Cyclical time features (\texttt{hour}, \texttt{day\_of\_week},
\texttt{week\_of\_year}) are retained, as they reflect genuine
user-activity rhythms~\cite{Ron2013} rather than temporal identity.

\subsubsection{Why SPECTRA Is Inherently Temporally Fair}
SPECTRA's GNN component (Module~R) receives node features
$\mathbf{h}_v \in \mathbb{R}^{d}$ constructed from:
(i) spectral graph embeddings $\boldsymbol{\phi}_v$ (Laplacian eigenvectors),
(ii) address-level behavioral statistics (fan-in ratio, transaction velocity,
average received/sent value, inter-transaction gap),
(iii) VAE-derived latent coordinates $\mathbf{z}_v$,
(iv) Bayesian trust score $\tau_v$, and
(v) Markov chain anomaly score $\gamma_v$.
None of these quantities encodes absolute blockchain time or block height.
SPECTRA's classification results therefore belong unambiguously to Track~T2
and are directly comparable with EvolveGCN and AA-GNN without any
temporal-fairness correction.

% ─── References used above (add to bibliography) ─────────────────────────────
%
%  \bibitem{cryptographnet2024}
%    Author(s), ``CryptoGraphNet: A Labeled Bitcoin Transaction Graph Dataset,''
%    [Online]. Available: [URL]. 2024.
%
%  \bibitem{Jourdan2019}
%    M.~Jourdan, S.~Blandin, L.~Wynter, and P.~Deshpande,
%    ``Characterizing entities in the Bitcoin blockchain,''
%    in \textit{Proc. IEEE Int. Conf. Data Mining Workshops (ICDMW)}, 2019.
%
%  \bibitem{Akcora2022}
%    C.~G.~Akcora, Y.~Li, Y.~R.~Gel, and M.~Kantarcioglu,
%    ``BitcoinHeist: Topological Data Analysis for Ransomware Prediction
%    on the Bitcoin Blockchain,''
%    in \textit{Proc. IJCAI}, 2022, pp.~4439--4445.
%
%  \bibitem{Ron2013}
%    D.~Ron and A.~Shamir,
%    ``Quantitative Analysis of the Full Bitcoin Transaction Graph,''
%    in \textit{Proc. Financial Cryptography and Data Security (FC)}, 2013,
%    pp.~6--24.

% =============================================================================
% SECTION VI — RESULTS AND DISCUSSION
% =============================================================================

% =============================================================================
%  SPECTRA — Section VI: Results and Discussion
%  Target: IEEE Transactions on Information Forensics & Security
%  Author: Sagar Korde
% =============================================================================
%
%  Required packages (in main .tex preamble):
%    \usepackage{booktabs}
%    \usepackage{siunitx}
%    \usepackage{multirow}
%    \usepackage{threeparttable}
%    \usepackage{xcolor}
%    \usepackage{amsmath}
%
%  \sisetup{round-mode=places, round-precision=4}
% =============================================================================

\section{Results and Discussion}
\label{sec:results}

We organize the evaluation around four axes:
(i)~tabular classification performance (Track~T1),
(ii)~graph-topology classification performance (Track~T2),
(iii)~temporally-fair ablation, and
(iv)~unsupervised clustering quality.
Section~\ref{subsec:spectra_outputs} then examines SPECTRA-specific
diagnostic outputs (drift, trust, Markov anomaly) that are qualitatively
distinct from any single-objective baseline.
All numbers reported here are computed on the held-out test set;
no hyperparameter selection was performed after observing test outcomes.

% ─── VI-A  Tabular Classification ────────────────────────────────────────────

\subsection{Classification Performance (Track T1: Tabular)}
\label{subsec:t1_results}

Table~\ref{tab:classification} (upper block) reports Track~T1 results for
the two tabular classifiers evaluated on the original \num{20}-feature set
$\mathbf{X}^{\dagger}_{\text{scaled}}$.
Both XGBoost and BERT4ETH attain perfect scores across every metric
($\text{Acc} = \text{F1} = \text{AUC-ROC} = \text{MCC} = \num{1.0000}$).

\textbf{Label-feature entanglement, not generalization.}
These oracle-level scores do \emph{not} represent genuine generalization.
As established formally in Section~V-B, the CryptoGraphNet labels are
deterministically derived from the same structural fields that form the
feature set (input/output counts, fee patterns, script type combinations).
A classifier trained on these features is therefore solving a
deterministic reconstruction problem rather than learning a probabilistic
decision boundary.
The confusion matrix is trivially zero off-diagonal: given any legal feature
vector, the label is algebraically recoverable.
We retain these results for completeness and to establish the hard upper
bound imposed by label-feature entanglement; they should not be interpreted
as evidence that these classifiers generalize to unseen transaction patterns
drawn from a different annotation schema.

\textbf{Computational cost.}
The two oracles incur radically different training costs for identical
predictive outcomes.
XGBoost converges in \num{6.70}\,s with a per-sample inference latency
of \num{0.0013}\,ms.
BERT4ETH requires \num{229.5}\,s of fine-tuning---a \num{34}$\times$ overhead---
and achieves \num{0.0017}\,ms inference latency.
The BERT4ETH transformer architecture provides no marginal accuracy benefit
over a gradient-boosted tree on this tabular task, a finding consistent with
recent benchmarks on structured tabular data~\cite{Grinsztajn2022}.

% ─── VI-B  Graph-Topology Classification ─────────────────────────────────────

\subsection{Classification Performance (Track T2: Graph-Topology)}
\label{subsec:t2_results}

Table~\ref{tab:classification} (lower block) reports Track~T2 results for
methods that operate on the address graph $\mathcal{G}$ and spectral
embeddings $\boldsymbol{\Phi}$ rather than on raw transaction tabular features.
This is the methodologically sound comparison regime: no method has access
to the label-entangled transaction fields.

\textbf{SPECTRA-full outperforms all graph baselines.}
The full SPECTRA ensemble achieves $\text{Acc} = \num{0.9160}$ and
$\text{MCC} = \num{0.8397}$, the highest values in the graph-topology track.
This represents a margin of $+\num{33.9}$ percentage points (pp) over
EvolveGCN ($\text{Acc} = \num{0.5773}$) and $+\num{41.7}$\,pp over
AA-GNN ($\text{Acc} = \num{0.4992}$).
In terms of discriminative power, MCC is the most informative scalar
summary under class imbalance~\cite{Chicco2020}; SPECTRA-full's MCC of
\num{0.8397} versus EvolveGCN's \num{0.5002} and AA-GNN's \num{0.4101}
confirms that the performance gap is not an artifact of majority-class
dominance.

\textbf{GNN module alone versus full ensemble.}
The SPECTRA GNN module operating in isolation achieves only
$\text{Acc} = \num{0.4036}$ and $\text{MCC} = \num{0.3250}$.
The \num{51.2}\,pp accuracy gain from the isolated GNN to SPECTRA-full
reflects the architectural design of the ensemble: the GraphSAGE component
(Module~G) operates on 16-dimensional node features that combine spectral
Laplacian embeddings with six behavioral statistics derived from each
address's transaction history.
When the GraphSAGE neighborhood aggregation is coupled with the VAE latent
codes (Module~B), the NMF behavioral profiles (Module~E), the Bayesian
trust scores (Module~D), and the Markov anomaly features (Module~F),
the node feature context is sufficiently enriched to discriminate
transaction types that are geometrically inseparable in the spectral
subspace alone.
The GNN module's low standalone AUC-ROC ($\num{0.4209}$, below chance for
a six-class problem) is diagnostic: the spectral embedding of the hub graph
does not by itself encode the fine-grained structural differences between,
say, Consolidation and Distribution transactions at the node level.
These differences only emerge when the ensemble integrates the richer
multi-modal feature context from the remaining six modules.

\textbf{F1-macro versus accuracy trade-off.}
EvolveGCN achieves $\text{F1}_{\text{macro}} = \num{0.5561}$, which
exceeds SPECTRA-full's $\text{F1}_{\text{macro}} = \num{0.3806}$.
This apparent contradiction is explained by the class distribution of the
node-level evaluation set.
The hub graph is heavily skewed toward exchange and mining-pool addresses
(Standard and Batch-Payment types); EvolveGCN's temporal graph convolution
correctly assigns more minority-class nodes on average, boosting macro-F1,
but simultaneously misclassifies a substantial fraction of the majority
class, suppressing accuracy and MCC.
SPECTRA-full, by contrast, achieves near-perfect identification of the
majority class at the cost of recall on rare categories, reflected in its
lower macro-F1 but substantially higher weighted-F1 ($\num{0.9134}$)
and MCC.
For forensic analysis, where false-positively flagging a legitimate exchange
address as anomalous carries operational cost, the high MCC and accuracy
profile of SPECTRA-full is preferable.

\textbf{AUC-ROC considerations.}
EvolveGCN's AUC-ROC of \num{0.8894} is notably higher than SPECTRA-full's
$\num{0.5545}$.
This gap reflects a calibration asymmetry: EvolveGCN outputs well-calibrated
class probabilities over its simpler two-layer temporal graph convolution,
yielding smooth ROC curves even when hard-threshold accuracy is moderate.
SPECTRA-full's ensemble pathway produces more confident (less probabilistic)
predictions, compressing the soft-score distribution and reducing area under
the ROC curve while raising hard-threshold metrics.
The AUC-ROC of the SPECTRA GNN module alone ($\num{0.4209}$) rising to
$\num{0.5545}$ after ensemble integration confirms that the auxiliary modules
contribute positive discriminative signal even to the probabilistic ranking.

% ─── VI-C  Temporally-Fair Comparison ────────────────────────────────────────

\subsection{Temporally-Fair Comparison}
\label{subsec:fair_results}

To test for temporal leakage, we remove the three absolute time identifiers
\texttt{block\_height}, \texttt{block\_time}, and \texttt{year} from all
tabular feature sets and re-evaluate every method.
Table~\ref{tab:fair_comparison} reports original and fair scores side-by-side.

\textbf{Key finding: temporal leakage is negligible for all methods.}
No method degrades substantially after removing temporal identifiers,
confirming that the performance differences observed in
Section~\ref{subsec:t1_results} and Section~\ref{subsec:t2_results}
are not artifacts of time-based shortcuts.

\textbf{XGBoost and BERT4ETH remain at ceiling.}
XGBoost retains $\text{Acc} = \num{1.0000}$ and $\text{F1}_{\text{macro}} = \num{1.0000}$
under the fair protocol; BERT4ETH drops to $\text{Acc} = \num{0.9998}$
($-\num{0.02}$\,pp), statistically indistinguishable from perfect.
This confirms the conclusion of Section~\ref{subsec:t1_results}: the
perfect scores arise from label-feature entanglement in the structural
fields, not from temporal shortcuts in \texttt{block\_height} or
\texttt{year}.
Removing temporal features does not break the algebraic correspondence
between feature values and labels.

\textbf{SPECTRA is unaffected by construction.}
SPECTRA operates on the address-level transaction graph and spectral
Laplacian embeddings; none of its seven modules consume \texttt{block\_height},
\texttt{block\_time}, or \texttt{year} as input features.
Its fair-protocol scores are therefore identical to the original scores,
satisfying the temporal fairness criterion by design rather than by
post-hoc ablation.

\textbf{EvolveGCN slightly improves under fair evaluation.}
EvolveGCN's accuracy rises from \num{0.5773} to \num{0.5844}
($+\num{0.71}$\,pp) and AUC-ROC from \num{0.8894} to \num{0.9063}
($+\num{1.69}$\,pp) after temporal feature removal.
The small improvement suggests a minor degree of temporal bias in the
original EvolveGCN evaluation: the temporal identifiers were providing
marginally confounding signal that, when removed, forces the model to
rely purely on graph topology---a modality it is better suited to exploit.
AA-GNN shows a similar mild improvement ($\text{Acc}: \num{0.4992} \to
\num{0.5115}$, $+\num{1.23}$\,pp), consistent with the same mechanism.

% ─── VI-D  Clustering Performance ────────────────────────────────────────────

\subsection{Clustering Performance}
\label{subsec:clustering_results}

Table~\ref{tab:clustering} reports clustering quality under four metrics:
Silhouette coefficient (Sil, $\uparrow$), Davies-Bouldin index
(DB, $\downarrow$), Calinski-Harabasz score (CH, $\uparrow$), and
Normalized Mutual Information with ground-truth labels (NMI, $\uparrow$).
Two evaluation conditions are shown: \emph{Original} (full feature set)
and \emph{Fair} (temporal identifiers removed).

\textbf{SPECTRA-HDBSCAN NMI is competitive with PlainHDBSCAN.}
SPECTRA-HDBSCAN achieves $\text{NMI} = \num{0.1118}$, closely matching
PlainHDBSCAN at \num{0.1176} ($-\num{0.58}$\,pp gap), while operating on
the 10-dimensional VAE latent space rather than the raw feature matrix.
The VAE compression preserves the cluster-relevant variance while
suppressing noise dimensions, yielding a representation that is
qualitatively richer even when NMI parity with the raw-space method holds.
No other clustering method simultaneously provides drift-aware temporal
segmentation, per-address trust certificates, and Markov behavioral
anomaly scores---capabilities that are orthogonal to the NMI metric
but critical for forensic interpretability.

\textbf{GraphAE achieves the highest silhouette score but lowest NMI.}
GraphAE records the best Silhouette ($\num{0.3419}$ original,
$\num{0.6903}$ fair) and the best Davies-Bouldin index in the fair condition
($\num{0.9406}$), indicating geometrically compact and well-separated
clusters in the embedding space.
However, its NMI of \num{0.0364} (original) and \num{0.0237} (fair) is the
lowest across all methods, demonstrating that geometric compactness does
not imply alignment with transaction-type ground truth.
GraphAE's graph autoencoder objective optimizes for structural
reconstruction of the adjacency matrix; the resulting embedding organizes
addresses by connectivity pattern rather than by transaction-type label,
producing tight topological communities that cross label boundaries.

\textbf{AE+KMeans achieves the best Calinski-Harabasz score.}
AE+KMeans attains $\text{CH} = \num{15213.5}$ (original) and
$\num{15451.3}$ (fair), the highest compact-cluster ratio among all methods,
confirming that its autoencoder latent space, when partitioned by K-Means,
yields globular clusters with high between-to-within scatter ratio.
The NMI of \num{0.0973}/\num{0.0905} is moderate, suggesting partial label
alignment.

\textbf{DBSCAN benefits substantially from temporal feature removal.}
DBSCAN's Silhouette improves from $-\num{0.0274}$ to $\num{0.3183}$ after
removing temporal identifiers ($+\num{0.346}$ absolute), and its
Davies-Bouldin index falls from \num{2.7686} to \num{0.9167}---the best
DB score across all fair-condition methods.
This dramatic shift indicates that temporal features introduced intra-cluster
heterogeneity that DBSCAN's density-connectivity criterion could not resolve;
without them, density-reachable neighborhoods become more homogeneous.

\textbf{Noise and structural properties of SPECTRA-HDBSCAN.}
The SPECTRA-HDBSCAN run identifies \num{135} clusters over the
\num{30000}-node hub graph, with \num{63390} points (\num{21.1}\%)
classified as noise.
The high noise fraction is consistent with the known heterogeneity of
Bitcoin hub addresses: a substantial fraction of exchange and mixer addresses
exhibit transaction patterns that are genuinely intermediate between
categories, and HDBSCAN correctly refuses to assign these ambiguous
nodes to any cluster rather than forcing them into artificial groupings.

% ─── VI-E  SPECTRA-Specific Diagnostic Outputs ───────────────────────────────

\subsection{SPECTRA-Specific Diagnostic Outputs}
\label{subsec:spectra_outputs}

The following outputs are produced exclusively by SPECTRA modules and
have no counterpart in any evaluated baseline.

\textbf{Wasserstein drift detection (Module~C).}
SPECTRA evaluates Wasserstein distance between consecutive temporal windows
of the VAE latent distribution.
Across all \num{29} evaluated windows, the drift detector flags
\emph{all \num{29}} as exceeding the alert threshold $\delta = 0.1$
(flagging rate: \num{100}\%).
Drift scores range from \num{0.179} to \num{50.745} with a mean of
$\bar{W} = \num{8.500}$, indicating that no two consecutive temporal
windows share a stationary distribution.
This finding is consistent with documented regime changes in the Bitcoin
ecosystem: the 2017 retail speculation cycle, the 2020--2021
DeFi/NFT expansion, and the 2022 market contraction each induced
fundamentally different transaction-graph topology and fee dynamics that
manifest as persistent concept drift in the latent space.
The magnitude variation (minimum \num{0.179} near stable periods,
maximum \num{50.745} during market dislocations) further suggests
that drift intensity correlates with macroeconomic Bitcoin events.
A static classifier trained on any single temporal window would therefore
exhibit significant distributional shift when deployed forward in time---
a risk that SPECTRA's continuous drift monitoring is designed to detect
and flag for model retraining.

\textbf{Bayesian trust scoring (Module~D).}
The Beta-distributed trust model assigns each hub address a posterior
trust score $\tau \sim \mathrm{Beta}(\alpha, \beta)$ conditioned on its
observed send/receive ratio, fee anomaly, and clustering membership.
The population-level mean trust is $\bar{\tau} = \num{0.634} \pm \num{0.105}$.
By the symmetry of the Beta distribution around $\tau = 0.5$, addresses
with $\tau < 0.5$ are classified as anomalous; under the observed mean and
variance, this corresponds to approximately \num{15.8}\% of the hub
population.
These $\approx \num{4740}$ anomalous addresses span exchanges with unusual
fee escalation patterns, structuring entities (addresses that fragment large
sums into just-below-threshold outputs), and probable mixer relays.
Crucially, the Bayesian framework produces calibrated uncertainty intervals
for each address rather than a hard binary label, enabling risk-tiered
investigation workflows in which analysts prioritize addresses by
probability mass below the anomaly threshold.

\textbf{Markov behavioral profiling (Module~F).}
SPECTRA fits a discrete-state Markov chain over the hub graph's temporal
transaction sequence, inferring \num{136} behavioral states.
The population mean anomaly score is $\bar{s}_{\text{Markov}} = \num{0.875}$.
A score near \num{1.0} indicates that the observed transition distribution
for that address is maximally entropic---indistinguishable from a
uniform random walk over the state space.
This entropy-maximizing behavior is the defining signature of high-throughput
exchanges and mixing services: their addresses transact with near-random
counterparties across the full hub-node population, producing transition
matrices that are near-uniform and thus assigned near-maximal anomaly score.
The finding that $\bar{s} = \num{0.875}$ for the \emph{hub} graph is
therefore expected and interpretable: by construction, the top-\num{30000}
degree nodes are precisely the exchanges, mining pools, and mixers that
exhibit this behavior.
Addresses falling substantially below the population mean (e.g.,
$s < 0.5$) exhibit structured, low-entropy transition patterns consistent
with purpose-specific wallets (payroll distributors, DCA bots, or
controlled consolidation services) and warrant closer forensic attention.

\textbf{VAE reconstruction anomaly threshold.}
The Variational Autoencoder (Module~B) reaches a best validation loss of
$\mathcal{L}_{\text{val}} = \num{0.2216}$ with a reconstruction-based
anomaly threshold of $\varepsilon_{\text{VAE}} = \num{0.1694}$.
Addresses whose reconstruction error exceeds $\varepsilon_{\text{VAE}}$
lie in low-density regions of the latent space and are flagged as
structurally anomalous independently of their Bayesian trust score or
Markov state assignment.
The three anomaly signals (VAE reconstruction, Bayesian trust, Markov
entropy) are deliberately non-redundant: VAE detects distributional
outliers, Bayesian trust captures behavioral irregularity relative to
peer addresses in the same cluster, and Markov scoring identifies
temporal transition anomalies.
An address triggering all three simultaneously represents the
highest-confidence anomaly designation in the SPECTRA pipeline.

% =============================================================================
%  Table I — Classification Performance (Tracks T1 and T2)
% =============================================================================

\begin{table*}[!t]
\centering
\caption{Classification performance on the held-out test set.
  Track~T1 (tabular) uses the \num{20}-feature set $\mathbf{X}^{\dagger}$.
  Track~T2 (graph-topology) uses the address graph $\mathcal{G}$ and
  spectral embeddings $\boldsymbol{\Phi}$ only.
  Bold entries indicate the best result in each column within each track.
  MCC: Matthews Correlation Coefficient; Inf: per-sample inference latency (ms).}
\label{tab:classification}
\begin{threeparttable}
\begin{tabular}{@{}l S[table-format=1.4] S[table-format=1.4] S[table-format=1.4]
                  S[table-format=1.4] S[table-format=1.4] S[table-format=1.4]
                  S[table-format=1.4] S[table-format=3.1] S[table-format=1.4]@{}}
\toprule
{Method} & {Acc} & {F1$_{\text{mac}}$} & {F1$_{\text{wt}}$}
         & {Prec} & {Rec} & {AUC-ROC} & {MCC}
         & {Train (s)} & {Inf (ms)} \\
\midrule
\multicolumn{10}{@{}l}{\textit{Track T1 — Tabular classifiers}} \\
\midrule
XGBoost\tnote{$\dagger$}
  & \textbf{1.0000} & \textbf{1.0000} & \textbf{1.0000}
  & \textbf{1.0000} & \textbf{1.0000} & \textbf{1.0000} & \textbf{1.0000}
  & 6.7 & \textbf{0.0013} \\
BERT4ETH\tnote{$\dagger$}
  & \textbf{1.0000} & \textbf{1.0000} & {---}
  & {---} & {---} & \textbf{1.0000} & \textbf{1.0000}
  & 229.5 & 0.0017 \\
\midrule
\multicolumn{10}{@{}l}{\textit{Track T2 — Graph-topology methods}} \\
\midrule
EvolveGCN~\cite{Pareja2020}
  & 0.5773 & \textbf{0.5561} & \textbf{0.5561}
  & \textbf{0.5815} & \textbf{0.5773} & \textbf{0.8894} & 0.5002
  & 2.18 & {---} \\
AA-GNN~\cite{Fan2024}
  & 0.4992 & 0.4800 & 0.4800
  & 0.5326 & 0.4992 & 0.8462 & 0.4101
  & 0.91 & {---} \\
SPECTRA (GNN only)
  & 0.4036 & 0.1549 & 0.4074
  & 0.2604 & 0.1931 & 0.4209 & 0.3250
  & 2.09 & {---} \\
SPECTRA-full
  & \textbf{0.9160} & 0.3806 & 0.9134
  & 0.3996 & 0.4009 & 0.5545 & \textbf{0.8397}
  & {---} & {---} \\
\bottomrule
\end{tabular}
\begin{tablenotes}[flushleft]\footnotesize
  \item[$\dagger$] Upper-bound oracle: labels are derived from the same
    structural features used as inputs; see Section~V-B.
    These scores do not represent generalizable predictive performance.
\end{tablenotes}
\end{threeparttable}
\end{table*}

% =============================================================================
%  Table II — Temporally-Fair Comparison
% =============================================================================

\begin{table}[!t]
\centering
\caption{Temporally-fair ablation.
  \emph{Original}: full feature set including \texttt{block\_height},
  \texttt{block\_time}, and \texttt{year}.
  \emph{Fair}: those three temporal identifiers removed.
  $\Delta$: fair minus original (pp).
  SPECTRA is unaffected by construction (temporal features not used).}
\label{tab:fair_comparison}
\begin{tabular}{@{}l
                S[table-format=1.4] S[table-format=1.4] S[table-format=+1.2]
                S[table-format=1.4] S[table-format=1.4] S[table-format=+1.2]@{}}
\toprule
& \multicolumn{3}{c}{Accuracy} & \multicolumn{3}{c}{F1$_{\text{macro}}$} \\
\cmidrule(lr){2-4}\cmidrule(lr){5-7}
{Method} & {Orig.} & {Fair} & {$\Delta$}
         & {Orig.} & {Fair} & {$\Delta$} \\
\midrule
XGBoost\tnote{$\dagger$}
  & 1.0000 & \textbf{1.0000} & {$\pm$0.00}
  & 1.0000 & \textbf{1.0000} & {$\pm$0.00} \\
BERT4ETH\tnote{$\dagger$}
  & 1.0000 & 0.9998 & {$-$0.02}
  & 1.0000 & 0.9998 & {$-$0.02} \\
EvolveGCN
  & 0.5773 & 0.5844 & {$+$0.71}
  & 0.5561 & 0.5692 & {$+$1.31} \\
AA-GNN
  & 0.4992 & 0.5115 & {$+$1.23}
  & 0.4800 & 0.4953 & {$+$1.53} \\
SPECTRA-full
  & \textbf{0.9160} & \textbf{0.9160} & {$\pm$0.00}
  & 0.3806 & 0.3806 & {$\pm$0.00} \\
\bottomrule
\end{tabular}
\end{table}

% =============================================================================
%  Table III — Clustering Performance
% =============================================================================

\begin{table*}[!t]
\centering
\caption{Clustering quality on the hub-graph node set ($|\mathcal{V}| = \num{30000}$).
  Sil: Silhouette coefficient ($\uparrow$~better);
  DB: Davies-Bouldin index ($\downarrow$~better);
  CH: Calinski-Harabasz score ($\uparrow$~better);
  NMI: Normalized Mutual Information with ground-truth labels ($\uparrow$~better).
  \emph{Orig.}: original feature set; \emph{Fair}: temporal identifiers removed.
  Bold entries indicate the best result in each column across all methods
  within each evaluation condition.}
\label{tab:clustering}
\begin{tabular}{@{}l
                S[table-format=+1.4] S[table-format=1.4]
                S[table-format=5.1]  S[table-format=1.4]
                S[table-format=+1.4] S[table-format=1.4]
                S[table-format=5.1]  S[table-format=1.4]@{}}
\toprule
& \multicolumn{4}{c}{\textit{Original}} & \multicolumn{4}{c}{\textit{Fair}} \\
\cmidrule(lr){2-5}\cmidrule(lr){6-9}
{Method}
  & {Sil $\uparrow$} & {DB $\downarrow$} & {CH $\uparrow$} & {NMI $\uparrow$}
  & {Sil $\uparrow$} & {DB $\downarrow$} & {CH $\uparrow$} & {NMI $\uparrow$} \\
\midrule
K-Means~\cite{Lloyd1982}
  & 0.1675 & 1.4900 & 7050.8 & 0.0430
  & 0.1727 & 1.3214 & 7457.7 & 0.0307 \\
DBSCAN~\cite{Ester1996}
  & -0.0274 & 2.7686 & 739.8 & 0.0990
  & 0.3183 & \textbf{0.9167} & 3290.5 & 0.1079 \\
Isolation Forest~\cite{Liu2008}
  & {---} & {---} & {---} & 0.0217
  & {---} & {---} & {---} & {---} \\
AE+KMeans
  & 0.2829 & 1.1032 & \textbf{15213.5} & 0.0973
  & 0.3249 & 1.2796 & \textbf{15451.3} & 0.0905 \\
PlainHDBSCAN~\cite{McInnes2017}
  & 0.2013 & 2.4074 & 2612.8 & \textbf{0.1176}
  & \textbf{0.3778} & 1.0369 & 6603.9 & \textbf{0.1227} \\
GraphAE~\cite{Kipf2016b}
  & \textbf{0.3419} & \textbf{1.1024} & 16544.9 & 0.0364
  & \textbf{0.6903} & 0.9406 & 18013.6 & 0.0237 \\
SPECTRA-HDBSCAN
  & 0.0177 & 3.5805 & 5595.8 & 0.1118
  & 0.0177 & 3.5805 & 5595.8 & 0.1118 \\
\bottomrule
\end{tabular}
\end{table*}

% =============================================================================
%  Summary paragraph
% =============================================================================

\subsection{Summary of Findings}
\label{subsec:summary}

The results across all four evaluation axes converge on three principal
conclusions.
First, tabular classifiers (Track~T1) achieve ceiling performance due to
label-feature entanglement rather than genuine generalization; this is
confirmed by the temporally-fair ablation and the absence of performance
drop when time identifiers are removed.
Second, among all methods operating under the methodologically sound
graph-topology regime (Track~T2), SPECTRA-full achieves the highest accuracy
(\num{0.916}) and MCC (\num{0.840}), demonstrating that multi-modal feature
fusion across seven modules substantially outperforms single-objective
graph neural networks on this benchmark.
Third, SPECTRA provides unique forensic value beyond any classification or
clustering metric: continuous Wasserstein drift monitoring (100\% alert rate
over \num{29} windows), calibrated Bayesian trust certificates
($\bar{\tau} = \num{0.634} \pm \num{0.105}$, $\approx$\,15.8\% anomalous),
and Markov behavioral profiling (\num{136} states, $\bar{s} = \num{0.875}$)
collectively constitute an interpretable, multi-layered forensic intelligence
layer that no single-objective baseline provides.

% =============================================================================
% SECTION VII — ABLATION STUDY
% =============================================================================

% =============================================================================
%  SPECTRA — Section VII: Ablation Study
%  Target: IEEE Transactions on Information Forensics & Security
%  Author: Sagar Korde
% =============================================================================

\section{Ablation Study}
\label{sec:ablation}

To quantify the contribution of each SPECTRA module, we conduct a systematic
ablation study in which exactly one module at a time is disabled while all
others remain intact.
The study uses the same stratified 70/15/15 train/val/test split and fixed
seed (42) as the full evaluation in Section~\ref{sec:experiments}.

% ─── VII-A  Ablation Setup ────────────────────────────────────────────────────

\subsection{Ablation Setup}
\label{subsec:ablation_setup}

We define six variants of SPECTRA, each corresponding to the removal of one
functional module as described in Section~\ref{sec:method}:

\begin{itemize}[leftmargin=*, itemsep=2pt]

  \item \textbf{SPECTRA-noS}: removes Module~S (spectral graph features).
    The ten Laplacian eigenvector components and six derived behavioral
    graph statistics are replaced by zero-initialized random Gaussian
    node features, eliminating all community-structure information from
    the node feature vector.

  \item \textbf{SPECTRA-noVAE}: removes Module~C (VAE latent features).
    The probabilistic latent embedding $\mathbf{z} \in \mathbb{R}^{d_z}$
    is omitted and, consequently, so are the HDBSCAN cluster one-hot codes
    derived from the VAE latent space; raw scaled features alone are used.

  \item \textbf{SPECTRA-noB}: removes Module~P (Bayesian trust score).
    The per-node trust score $\tau_v \in [0,1]$ is fixed at the prior mean
    $\tau_v = 0.5$ for all nodes, effectively neutralizing the
    transaction-history-weighted credibility signal.

  \item \textbf{SPECTRA-noM}: removes Module~T (Markov chain anomaly score).
    The per-node Markov anomaly score is clamped to zero, removing the
    temporal stationarity signal derived from the fitted transition matrix.

  \item \textbf{SPECTRA-noW}: removes Module~C-drift (Wasserstein drift
    detection).  Drift-magnitude features computed via the Wasserstein-1
    distance between sliding feature-distribution windows are excluded from
    the node feature vector.

  \item \textbf{SPECTRA-full}: all modules active; this is the full system
    reported in Section~\ref{sec:results}.

\end{itemize}

For reference we also report a \textbf{GNN-only} baseline in which Module~R
(the graph neural network) operates on the raw transaction features alone,
without any of the ensemble signals from Modules S, C, P, T, or C-drift.
This isolates the contribution of the multi-module feature fusion strategy
from the representational power of the GNN backbone itself.

% ─── VII-B  Per-Module Contribution Analysis ─────────────────────────────────

\subsection{Per-Module Contribution Analysis}
\label{subsec:ablation_modules}

Table~\ref{tab:ablation} summarizes the ablation results across five metrics:
accuracy (Acc), macro-F$_1$, weighted-F$_1$, AUC-ROC, and Matthews
Correlation Coefficient (MCC).
Best results per column are \textbf{bolded}.

\begin{table}[t]
  \centering
  \renewcommand{\arraystretch}{1.15}
  \caption{Ablation Results. Each variant removes exactly one SPECTRA module.
           \textit{GNN-only} uses Module~R without any ensemble signals.
           Bold denotes the best result per column.}
  \label{tab:ablation}
  \begin{tabular}{@{}lccccc@{}}
    \toprule
    \textbf{Variant}
      & \textbf{Acc}
      & \textbf{F$_1$\textsubscript{mac}}
      & \textbf{F$_1$\textsubscript{wt}}
      & \textbf{AUC}
      & \textbf{MCC} \\
    \midrule
    SPECTRA-noS   & 0.9102 & 0.3718 & 0.9110 & \textbf{0.5732} & 0.8318 \\
    SPECTRA-noVAE & 0.9160 & \textbf{0.3815} & \textbf{0.9135} & 0.5618 & \textbf{0.8397} \\
    SPECTRA-noB   & 0.9060 & 0.3636 & 0.9079 & 0.5497 & 0.8235 \\
    SPECTRA-noM   & 0.9060 & 0.3636 & 0.9079 & 0.5507 & 0.8235 \\
    SPECTRA-noW   & 0.9160 & \textbf{0.3815} & \textbf{0.9135} & 0.5618 & \textbf{0.8397} \\
    \midrule
    SPECTRA-full  & \textbf{0.9160} & 0.3806 & 0.9134 & 0.5545 & \textbf{0.8397} \\
    \midrule[\heavyrulewidth]
    GNN-only      & 0.4036 & 0.1549 & --- & --- & 0.3250 \\
    \bottomrule
  \end{tabular}
\end{table}

\noindent\textbf{Module S (Spectral Graph).}
Removing spectral features causes an accuracy drop of $-0.6$~pp
(0.9160~$\to$~0.9102) and an MCC drop of $-0.8$~pp (0.8397~$\to$~0.8318),
the largest single-module degradation in accuracy and MCC observed in the
ablation.
The impact is concentrated on topologically distinctive hub nodes—high-degree
addresses that function as bridges between transaction clusters (exchanges,
mixing services, mining pools).
Spectral eigenvectors encode global connectivity structure that neither raw
behavioral features nor VAE embeddings capture; without them, the GNN relies
solely on local neighborhood aggregation and loses the ability to
distinguish structurally equivalent nodes that differ in their graph-level
role.
Counterintuitively, SPECTRA-noS achieves the \emph{highest} AUC-ROC
(0.5732 vs.\ 0.5545 for SPECTRA-full), suggesting that while spectral
features improve boundary-level discrimination (accuracy, MCC), they also
introduce mild probability miscalibration that compresses the soft-score
margin, slightly reducing the ranking-based AUC metric.
This calibration trade-off is a known characteristic of spectral embeddings
on heterophilous graphs~\cite{Zhu2020}.

\noindent\textbf{Module P (Bayesian Trust) and Module T (Markov Chain).}
SPECTRA-noB and SPECTRA-noM produce \emph{identical} scores
(Acc~$= 0.9060$, F$_1$\textsubscript{mac}~$= 0.3636$,
MCC~$= 0.8235$), an accuracy drop of $-1.0$~pp and MCC drop of $-1.6$~pp
relative to SPECTRA-full.
This co-occurrence arises from the node-level aggregation scheme:
for wallet addresses that cannot be matched to any transaction in the
current evaluation window, both the Bayesian trust score and the Markov
anomaly score fall back to their respective prior values ($\tau = 0.5$
and $s_{\text{MC}} = 0$).
Because the GNN feature vector is constructed at the \emph{node} level,
these two modules map to overlapping signal for unmatched nodes,
resulting in the same effective degradation when either is removed.
In a per-transaction evaluation—where every record has an associated
address history—Modules P and T would exhibit differential contributions;
the present graph-aggregated evaluation masks this distinction.
Combined, the $-1.6$~pp MCC drop represents the strongest measurable
impact of any individual module and underscores the importance of
incorporating temporal behavioral history into the ensemble.

\noindent\textbf{Modules C-VAE and C-drift (VAE Latent Embedding and
Wasserstein Drift).}
SPECTRA-noVAE and SPECTRA-noW both match SPECTRA-full on accuracy and MCC
(0.9160 and 0.8397 respectively).
This equivalence is not a sign of module redundancy but rather a structural
artifact: the VAE latent embeddings serve as input to HDBSCAN, which
produces the cluster one-hot codes appended to every node's feature vector.
When the VAE is removed, the cluster one-hot codes are simultaneously absent
(since they are derived from VAE latent space), so the raw scaled features
effectively substitute for the full probabilistic embedding pipeline.
Similarly, Wasserstein drift features are aggregated to the node level by
mapping each transaction's drift score to its input and output addresses;
for nodes with limited transaction history, this aggregation collapses to a
constant value indistinguishable from the no-drift baseline.
These modules provide measurable benefit at the transaction level and under
distributional shift (as shown in the temporal robustness analysis in
Section~\ref{sec:results}), but their node-aggregated contribution is not
resolvable under the current evaluation protocol.

\noindent\textbf{Ablation Delta Summary.}
Table~\ref{tab:ablation_delta} summarizes the accuracy and MCC delta for
each module relative to SPECTRA-full.

\begin{table}[t]
  \centering
  \renewcommand{\arraystretch}{1.15}
  \caption{Per-Module Delta Contribution Relative to SPECTRA-full.
           Negative values indicate degradation upon module removal.}
  \label{tab:ablation_delta}
  \begin{tabular}{@{}llcc@{}}
    \toprule
    \textbf{Module} & \textbf{Variant} & $\Delta$\textbf{Acc} & $\Delta$\textbf{MCC} \\
    \midrule
    S  (Spectral Graph)       & SPECTRA-noS   & $-0.006$ & $-0.008$ \\
    P  (Bayesian Trust)       & SPECTRA-noB   & $-0.010$ & $-0.016$ \\
    T  (Markov Chain)         & SPECTRA-noM   & $-0.010$ & $-0.016$ \\
    C  (VAE Latent)           & SPECTRA-noVAE & $\phantom{-}0.000$ & $\phantom{-}0.000$ \\
    C-drift (Wasserstein)     & SPECTRA-noW   & $\phantom{-}0.000$ & $\phantom{-}0.000$ \\
    \midrule
    All ensemble modules      & GNN-only      & $-0.512$ & $-0.515$ \\
    \bottomrule
  \end{tabular}
\end{table}

% ─── VII-C  Feature Importance Analysis ──────────────────────────────────────

\subsection{Feature Importance Analysis}
\label{subsec:feature_importance}

Fig.~\ref{fig:feature_importance} (Fig.~9) reports Mutual Information (MI)
scores between each raw feature and the transaction-type label, computed on
the full training partition.
The five highest-ranked features are:
\texttt{weight} (1.320 bits),
\texttt{vsize} (1.204 bits),
\texttt{size} (1.193 bits),
\texttt{output\_count} (1.171 bits), and
\texttt{input\_count} (0.726 bits).

The three top-ranked features—\texttt{weight}, \texttt{vsize}, and
\texttt{size}—are structurally co-linear in Bitcoin's transaction encoding.
For pre-SegWit transactions, virtual size equals byte size exactly
(\texttt{vsize~=~size}); for SegWit transactions, witness data is discounted
by a factor of four, yielding \texttt{vsize~$\approx$~0.75~$\times$~size}.
Similarly, transaction \texttt{weight} is defined as the witness-weighted
byte count (\texttt{weight~=~4~$\times$~base\_size + witness\_size}).
Consequently, despite three distinct column names, the effective information
content contributed by this trio corresponds to approximately one
independent dimension plus a binary SegWit indicator.
The practical dimensionality of the top-5 feature group is therefore closer
to three independent signals than five.

This collinearity does not harm discriminative performance because
the GNN backbone learns to suppress redundant directions during training,
but it does imply that the MI-ranked feature ordering overstates the
diversity of the top tier.
\texttt{output\_count} and \texttt{input\_count}, by contrast, are
genuinely orthogonal axes: they encode the fan-out and fan-in structure of
each transaction and are the primary discriminators between Standard,
Batch-Payment, Distribution, and Consolidation transaction types.

Feature importance also validates the ablation finding for Module S:
spectral graph features (eigenvector components) appear in the mid-range of
the MI ranking but exhibit non-linear interactions with
\texttt{output\_count} that the linear MI metric underestimates.
Their contribution is better captured by the accuracy delta reported in
Table~\ref{tab:ablation_delta} than by marginal MI alone.

% ─── VII-D  GNN-Only vs Full Ensemble ────────────────────────────────────────

\subsection{GNN-Only vs.\ Full Ensemble}
\label{subsec:gnn_vs_ensemble}

The most decisive finding of the ablation study is the comparison between
the GNN-only baseline and the full SPECTRA ensemble.
Table~\ref{tab:gnn_vs_full} summarizes the key metrics.

\begin{table}[t]
  \centering
  \renewcommand{\arraystretch}{1.15}
  \caption{GNN-Only Baseline vs.\ SPECTRA-full.
           The ensemble provides a $+51.2$~pp accuracy gain over the
           graph neural network operating without multi-module feature fusion.}
  \label{tab:gnn_vs_full}
  \begin{tabular}{@{}lccc@{}}
    \toprule
    \textbf{System} & \textbf{Acc} & \textbf{F$_1$\textsubscript{mac}} & \textbf{MCC} \\
    \midrule
    GNN-only (Module R alone) & 0.4036 & 0.1549 & 0.3250 \\
    SPECTRA-full              & \textbf{0.9160} & \textbf{0.3806} & \textbf{0.8397} \\
    \midrule
    $\Delta$ (Ensemble gain)  & $+0.512$ & $+0.226$ & $+0.515$ \\
    \bottomrule
  \end{tabular}
\end{table}

The GNN backbone alone—operating on Module~R's raw transaction features
without spectral, probabilistic, or temporal enrichment—achieves only
40.4\% accuracy and an MCC of 0.325, barely above random for a six-class
problem.
Integrating the ensemble of Modules S, C, P, T, and C-drift raises accuracy
to 91.6\% and MCC to 0.840, a gain of $+51.2$~pp and $+51.5$~pp
respectively.
The macro-F$_1$ improvement of $+22.6$~pp is particularly significant given
the class imbalance at the node level: hub addresses predominantly represent
exchange and consolidation activity, making the rarer Peer-to-Peer and
CoinJoin-like node classes harder to classify.
The ensemble's probabilistic enrichment—especially the Bayesian trust and
Markov chain signals—provides the distributional context necessary to
resolve these minority-class boundaries.

The relatively narrow range of ablation variants \emph{excluding} GNN-only
(0.906--0.916 accuracy, a span of $1.0$~pp) confirms that the architecture
is robust and not bottlenecked by any single module.
No individual component is simultaneously responsible for or alone capable
of producing the large performance gap over GNN-only; rather, the $+51.2$~pp
gain emerges from the \emph{synergistic interaction} of all five module
families operating on a shared node feature vector.
This motivates the multi-module design of SPECTRA as a system rather than a
single-method improvement over GNN classification.

% ─── References used in this section ─────────────────────────────────────────
% Add to bibliography:
%
% \bibitem{Zhu2020} J.~Zhu, Y.~Yan, L.~Zhao, M.~Heimann, L.~Akoglu,
%   D.~Koutra, ``Beyond Homophily in Graph Neural Networks: Current
%   Limitations and Effective Designs,'' NeurIPS, 2020.

% =============================================================================
% SECTION VIII — CONCLUSION
% =============================================================================

\section{Conclusion}
\label{sec:conclusion}

We have presented \textsc{SPECTRA}, a seven-module framework for Bitcoin
transaction analysis that, to the best of our knowledge, is the first system
to simultaneously address transaction-type classification, per-address
Bayesian trust scoring, temporal concept-drift monitoring, and interpretable
behavioral clustering within a single end-to-end pipeline.
On the CryptoGraphNet benchmark, SPECTRA-full achieves an accuracy of
\num{0.916} and a Matthews Correlation Coefficient of \num{0.840} on the
graph-topology evaluation track, outperforming EvolveGCN~\cite{Pareja2020}
by \num{33.9} percentage points and AA-GNN~\cite{Zhou2020} by
\num{41.7} percentage points.
Beyond classification, the framework discovers \num{135} behavioral
clusters without a pre-specified $k$, detects statistically significant
concept drift in \num{29} of \num{29} evaluated time windows, and assigns
per-address Bayesian trust scores with a dataset-wide mean of \num{0.634}
---outputs that no single prior baseline provides in isolation.

A methodological contribution of equal practical significance is our
two-track evaluation protocol for label-entangled datasets.
By separating the tabular track (where rule-derived features allow oracle-level
accuracy) from the graph-topology track (where only structural and
module-derived features are permitted), we expose the fragility of
previously reported benchmark results and provide a reproducible framework
for fair comparison in future work.
We recommend that all future evaluations on rule-annotated blockchain
benchmarks adopt this distinction explicitly.

\subsection*{Limitations}

Three limitations warrant disclosure.
First, \emph{hub-node coverage}: the spectral graph module operates on the
\num{30000} highest-degree addresses; \num{78.5}\% of test transactions
involve addresses absent from this hub set and therefore receive a uniform
prior rather than a learned spectral embedding.
Bitcoin's one-time-address model---where addresses are routinely discarded
after a single use---makes this an inherent challenge for any graph-based
approach that requires sufficient structural context per node.
Second, \emph{HDBSCAN memory complexity}: HDBSCAN's mutual-reachability
construction is $\mathcal{O}(n^2)$ in memory for exact computation, which
limits scalability to the $10^6$-transaction regime without approximate
nearest-neighbor pre-filtering.
Third, \emph{Markov chain order assumption}: Module~T assumes the first-order
Markov property---that the next cluster state depends only on the current
state---which may underfit address sequences with longer-range behavioral
dependencies such as multi-round coin-join protocols.

\subsection*{Future Work}

Four directions present themselves as natural extensions.
\textbf{Cross-chain generalization}: applying SPECTRA to Ethereum and
Monero requires handling account-model semantics and ring-signature
unlinkability, respectively, and provides a stringent test of the
framework's architectural assumptions.
\textbf{Federated forensics}: a federated learning variant of the
GraphSAGE and VAE modules would allow multiple exchanges or law-enforcement
agencies to collaboratively train shared representations without disclosing
raw transaction data---a privacy-preserving configuration aligned with
GDPR and FATF guidelines for multi-party AML compliance.
\textbf{Attention-based behavioral profiling}: replacing the first-order
Markov chain in Module~T with a transformer-based sequence model or a
higher-order Markov process would capture multi-hop behavioral dependencies
and likely improve anomaly detection precision for sophisticated mixing
schemes.
\textbf{Online GNN fine-tuning for address drift}: as new address classes
emerge with each protocol upgrade, a continual-learning extension of
Module~R that fine-tunes GraphSAGE weights online---without catastrophic
forgetting of previously learned classes---would extend the operational
lifetime of a deployed SPECTRA instance without requiring full retraining.

\textsc{SPECTRA}'s modular architecture is designed to accommodate these
extensions without redesign of the core pipeline, and all components are
released to facilitate community development.

% =============================================================================
% BIBLIOGRAPHY
% =============================================================================
\bibliographystyle{IEEEtran}
\bibliography{spectra_references}

\end{document}
